% Hadamard / Mittag-Leffler reduction for prob:weighted-rg-locality.
%
% This fragment proves the analytic part of the
% Costello-locality compatibility for the Hadamard parametrix on the
% Mittag-Leffler coefficient module of T1. The proof exploits the
% factorization (eq:wt-propagator-factor) of the BV propagator into a
% bulk spacetime piece and a coefficient piece: the spacetime piece
% carries all analytic content (Hadamard parametrix, Hörmander
% wave-front analysis), the coefficient piece is purely algebraic
% (the Casimir on the weighted Tate pair). The Mittag-Leffler limit
% in the coefficient direction therefore commutes with the
% bulk-spacetime locality argument, and locality of each windowwise
% truncation passes to the limit.
%
% The output is a coefficient/Hadamard criterion. The analytic
% wave-front and Mittag-Leffler parts of prob:weighted-rg-locality are
% proved graph by graph; the endpoint QME class remains the exact
% obstruction named in the theorem statements.

\subsection*{Hadamard parametrix on the Mittag-Leffler coefficient module}

\begin{rmk}[Spacetime/coefficient factorization of the Hadamard parametrix]\label{rmk:hadamard-factorization}
  In Costello's heat-kernel BV calculus
  \cite{costello-renormalization}*{Ch.~5}, the Hadamard parametrix
  enters as the asymptotic expansion of the heat kernel
  $K_t^{\mathrm{base}}$ on the bulk free elliptic complex
  $\Omega^\bullet(\R^2)\widehat\otimes\Omega^{0,\bullet}(\C^2)$,
  combined with Hörmander wave-front-set tracking
  \cite{hormander-vol1}*{Ch.~7} to control the singular support of the
  propagator $P_{\epsilon,L}^{\mathrm{base}}$ along the diagonal.
  Costello's locality package
  \cite{costello-renormalization}*{Chs.~5--7} is used here only as the
  heat-kernel/RG vocabulary for free elliptic complexes: local
  functionals admit Hadamard-parametrix expansions whose counterterms
  stay in the local-functional class $\mathcal O_{\mathrm{loc}}$ under
  the hypotheses of that formalism.  The exact published theorem number
  is not used as a local input.

  In the weighted setting of \ref{eq:wt-propagator-factor}, the
  propagator factors as
  $P_{\epsilon,L}^{w}=P_{\epsilon,L}^{\mathrm{base}}\widehat\otimes
   C_{\mathfrak h\oplus\mathfrak h^\vee[1],w}$. The bulk piece
  $P_{\epsilon,L}^{\mathrm{base}}$ is independent of the coefficient
  weight; the coefficient piece $C_{\mathfrak h,w}$ is purely
  algebraic, with no scale parameter and no Laplacian. The Hadamard
  parametrix construction therefore lives entirely on the bulk
  factor; the coefficient factor enters only by tensor.
\end{rmk}

\begin{lem}[Windowwise parametrix compatibility]\label{lem:windowwise-parametrix}
  For each $w_0\in\Z_{\geq1}$, let
  \[
    C_{\mathfrak h,w}^{w_0}
      =
    \sum_{d\leq w_0}\sum_i
      (w(d)e_{d,i})\otimes(w(d)^{-1}e^i_d)
      \in
    \Bigl(\prod_{d\leq w_0}w(d)H_d\Bigr)
      \otimes
    \Bigl(\prod_{d\leq w_0}w(d)^{-1}H_d^\vee\Bigr)
  \]
  denote the truncated weighted Casimir. Define the windowwise field complex
  \begin{equation}\label{eq:windowwise-field-complex}
    \mathcal E_w^{w_0}=\Omega^\bullet(\R^2)\widehat\otimes
      \Omega^{0,\bullet}(\C^2)\widehat\otimes
      \biggl(\Bigl(\prod_{d\leq w_0}w(d)H_d\Bigr)[1]
      \oplus\Bigl(\prod_{d\leq w_0}w(d)^{-1}H_d^\vee\Bigr)[2]\biggr)
  \end{equation}
  and the windowwise propagator
  $P_{\epsilon,L}^{w,w_0}=P_{\epsilon,L}^{\mathrm{base}}\widehat\otimes
  C_{\mathfrak h\oplus\mathfrak h^\vee[1],w}^{w_0}$. Then:
  \begin{enumerate}[label=(\arabic*)]
    \item For each $w_0$, $\mathcal E_w^{w_0}$ is a free elliptic
      complex with finite-rank coefficient bundle, and Costello's
      Hadamard parametrix construction
      \cite{costello-renormalization}*{Ch.~5} applies verbatim to
      $P_{\epsilon,L}^{w,w_0}$.
    \item The truncations $\{P_{\epsilon,L}^{w,w_0}\}_{w_0\geq1}$
      form a strict Mittag-Leffler inverse system in the coefficient
      direction, with surjective transition maps
      $P_{\epsilon,L}^{w,w_0'}\twoheadrightarrow P_{\epsilon,L}^{w,w_0}$
      for $w_0'\geq w_0$ given by truncation to degrees
      $d\leq w_0$.
    \item The wave-front set
      $\mathrm{WF}(P_{\epsilon,L}^{w,w_0})$ is contained in the
      bulk wave-front set $\mathrm{WF}(P_{\epsilon,L}^{\mathrm{base}})
      \subset T^*(\R^2\times\C^2)\setminus 0$, independently of
      $w_0$, since the coefficient factor
      $C_{\mathfrak h,w}^{w_0}$ is a finite-rank algebraic tensor
      with no spacetime singularities.
  \end{enumerate}
\end{lem}

\begin{proof}
  (1) For fixed $w_0$, the coefficient bundle
  $\bigl(\prod_{d\leq w_0}w(d)H_d\bigr)[1]\oplus
  \bigl(\prod_{d\leq w_0}w(d)^{-1}H_d^\vee\bigr)[2]$ is a finite-rank
  $\Z$-graded vector space (finite product of finite-dimensional
  $H_d$'s). Costello's Hadamard parametrix construction
  \cite{costello-renormalization}*{Ch.~5} applies verbatim to any free
  elliptic complex with finite-rank coefficient bundle, so it applies
  to $\mathcal E_w^{w_0}$. The bulk propagator
  $P_{\epsilon,L}^{\mathrm{base}}$ is unchanged; the coefficient
  factor $C_{\mathfrak h,w}^{w_0}$ is a fixed finite-rank tensor.

  (2) Surjectivity of the transition maps is the truncation
  $\prod_{d\leq w_0'}w(d)H_d\twoheadrightarrow\prod_{d\leq w_0}w(d)H_d$
  on the coefficient direction; the bulk factor is unchanged. An
  inverse system with surjective transition maps satisfies the
  Mittag-Leffler condition trivially, so the inverse limit is exact in
  the strict finite-window habitat of
  Definition~\ref{defn:strict-ml-tate-habitat}, and
  Proposition~\ref{prop:strict-ml-habitat-exactness} gives
  \(\varprojlim^1=0\) for this tower. Compatibility with
  $P_{\epsilon,L}^{w,w_0'}\mapsto P_{\epsilon,L}^{w,w_0}$ follows from
  the tensor factorization and the truncation of the Casimir.

  (3) The Hörmander wave-front set of a tensor product distribution
  $T_1\otimes T_2$ on $X_1\times X_2$ satisfies
  $\mathrm{WF}(T_1\otimes T_2)\subseteq\bigl(\mathrm{WF}(T_1)\times
  \{(x_2,0):x_2\in X_2\}\bigr)\cup\bigl(\{(x_1,0):x_1\in X_1\}\times
  \mathrm{WF}(T_2)\bigr)\cup\bigl(\mathrm{WF}(T_1)\times
  \mathrm{WF}(T_2)\bigr)$ (\cite{hormander-vol1}*{Thm.~8.2.9}). Here
  $T_1=P_{\epsilon,L}^{\mathrm{base}}$ has wave-front set in the
  conormal of the diagonal of $\R^2\times\C^2$, and
  $T_2=C_{\mathfrak h,w}^{w_0}$ is a finite-rank algebraic tensor
  with no point dependence, hence smooth, so
  $\mathrm{WF}(T_2)=\emptyset$. Hence
  $\mathrm{WF}(P_{\epsilon,L}^{w,w_0})=
  \mathrm{WF}(P_{\epsilon,L}^{\mathrm{base}})\times\{0\}$, which we
  identify with $\mathrm{WF}(P_{\epsilon,L}^{\mathrm{base}})$ on the
  spacetime, independent of $w_0$.
\end{proof}

\begin{thm}[Hadamard control on the Mittag-Leffler coefficient module]\label{thm:hadamard-mittag-leffler}
  Let $w(d)=q^d$ with $q>1$ be the exponential degree weight of
  Remark~\ref{rmk:wt-canonical-weight}, and let $\mathcal E_w$ be the
  weighted field complex \eqref{eq:wt-field-complex}. The Hadamard
  parametrix construction
  \cite{costello-renormalization}*{Ch.~5} extends from each
  finite-window restriction $\mathcal E_w^{w_0}$
  (Lemma~\ref{lem:windowwise-parametrix}) to the Mittag-Leffler
  limit $\mathcal E_w=\varprojlim_{w_0}\mathcal E_w^{w_0}$,
  preserving the Hadamard wave-front and finite-window locality
  estimates graph by graph. Specifically:
  \begin{enumerate}[label=(\roman*)]
    \item The propagator $P_{\epsilon,L}^{w}$ admits a Hadamard
      parametrix expansion as the inverse limit
      \[
        P_{\epsilon,L}^{w}
          =\varprojlim_{w_0}P_{\epsilon,L}^{w,w_0},
      \]
      with bulk wave-front set independent of $w_0$.
    \item For every local interaction
      $I\in\mathcal O_{\mathrm{loc}}(\mathcal E_w)$ of bounded
      polynomial-degree $D_0$ at each vertex in the coefficient
      direction, every fixed Costello graph $\Gamma$ of order
      $\hbar^n$ contributes a compatible finite-window counterterm
      system whose weighted coefficient size is bounded, on every
      finite output window, by a finite graph-dependent constant times
      $|C_n^{\mathrm{base}}[\Gamma]|$, where
      $|C_n^{\mathrm{base}}[\Gamma]|$ is the bulk spacetime
      counterterm of Costello.
    \item For every fixed graph whose finite-window counterterms are
      degreewise surjective and compatible under the Mittag-Leffler
      transition maps, the corresponding inverse-limit representative
      lies in the local-functional class
      \(\mathcal O_{\mathrm{loc}}(\mathcal E_w)\). The associated
      endpoint obstruction is the class in
      $H^1(\mathcal O_{\mathrm{loc}}(\mathcal E_w),Q+\{I_0,-\})$
      for the mixed brane-defect interaction.
    \item The construction has the ordinary strict inverse-limit
      universal property and no stronger one: a compatible graphwise
      counterterm family
      \(\{C_n[\Gamma]^{w_0}\}_{w_0}\) determines a unique element
      \(C_n[\Gamma]\in
      \varprojlim_{w_0}\mathcal O_{\mathrm{loc}}(\mathcal E_w^{w_0})\),
      and every finite-window projection of this element is the given
      counterterm.
  \end{enumerate}
  Thus the Hadamard and wave-front part of
  Problem~\ref{prob:weighted-rg-locality} is exactly the
  finite-window locality plus Mittag-Leffler compatibility statement
  above.  The theorem's hypotheses are graphwise; summability over graph
  topologies and QME vanishing are the separate obstruction data recorded
  in the displayed cohomology group.
\end{thm}

\begin{proof}
  (i) Each truncation $P_{\epsilon,L}^{w,w_0}$ admits Costello's
  Hadamard parametrix expansion by
  Lemma~\ref{lem:windowwise-parametrix}~(1). The Mittag-Leffler
  inverse system of (2) of that lemma is strict (surjective
  transition maps), so the inverse limit
  \[
    P_{\epsilon,L}^{w}
      =\varprojlim_{w_0}P_{\epsilon,L}^{w,w_0}
  \]
  exists in the weighted Tate category and inherits the wave-front
  bound (3) of that lemma uniformly in $w_0$. The inverse limit of the
  Hadamard expansions is itself a Hadamard expansion: the
  asymptotic Hadamard coefficients $a_k(x,y)$ on the bulk diagonal
  are independent of the coefficient direction, so the windowwise
  Hadamard expansion of $P_{\epsilon,L}^{w,w_0}$ is
  \[
    \bigl(\sum_{k\geq0}a_k(x,y)\,t^k\bigr)\widehat\otimes
      C_{\mathfrak h,w}^{w_0},
  \]
  and the inverse limit in the
  coefficient direction commutes with the asymptotic expansion in
  the bulk parameter $t$.

  (ii) Fix a graph $\Gamma$ with $V(\Gamma)$ vertices and
  $E(\Gamma)$ internal edges. Each vertex carries an interaction of
  polynomial degree $\leq D_0$ in the coefficient direction
  (hypothesis on $I$). Each internal edge contracts one source leg
  with one target leg via the Casimir $C_{\mathfrak h,w}$. By
  \ref{w:W1}, each such contraction multiplies the weight by a factor
  $\leq C_w$ on Hamiltonian legs, and the cotangent legs are bounded
  on the finite windows appearing in the graph by
  \ref{w:W2}. The total coefficient size after $E(\Gamma)$
  contractions is bounded by a finite graph-dependent constant times
  the bulk counterterm norm. The
  bulk counterterm norm $|C_n^{\mathrm{base}}[\Gamma]|$ is finite by
  Costello's bulk locality theorem
  \cite{costello-renormalization}*{Chs.~5--7}.

  This graph-dependent constant is finite for the fixed graph
  \(\Gamma\) and the fixed finite output window under consideration.
  A summability statement over all graph topologies at a fixed
  \(\hbar\)-order is not one of the hypotheses.  Such a statement belongs
  to the renormalized QME problem, together with the usual stability,
  counterterm, and locality hypotheses of Costello's graph expansion.
  The present theorem needs only the graphwise bound: for a fixed
  graph, the coefficient degrees and the corresponding weight ratios
  lie in a finite set, so the compatible finite-window counterterm
  system has a Mittag-Leffler limit when the windowwise representatives
  are themselves compatible.

  (iii) The bulk Costello locality theorem at fixed window $w_0$
  is used only as a finite-window input: the windowwise counterterm
  $C_n[\Gamma]^{w_0}\in\mathcal O_{\mathrm{loc}}(\mathcal E_w^{w_0})$
  is a local functional with kernel supported on the diagonal of
  $\R^2\times\C^2$. By the wave-front bound of
  Lemma~\ref{lem:windowwise-parametrix}~(3), the support is the
  bulk diagonal independently of $w_0$. The transition map
  $\mathcal O_{\mathrm{loc}}(\mathcal E_w^{w_0'})\to
  \mathcal O_{\mathrm{loc}}(\mathcal E_w^{w_0})$ for $w_0'\geq w_0$
  is the restriction of polynomial functionals to truncated
  coefficient indices, equivalently setting higher coefficient fields
  to zero; in the regulator-admissible sector these maps are
  degreewise surjective.  This preserves locality (the support of the
  truncated kernel is the same diagonal). The inverse limit
  $\mathcal O_{\mathrm{loc}}(\mathcal E_w)=
  \varprojlim_{w_0}\mathcal O_{\mathrm{loc}}(\mathcal E_w^{w_0})$
  exists in the weighted Tate category and is the candidate
  local-functional class on $\mathcal E_w$. The windowwise counterterms
  $\{C_n[\Gamma]^{w_0}\}_{w_0}$ form a compatible inverse system,
  and their inverse limit
  $C_n[\Gamma]=\varprojlim_{w_0}C_n[\Gamma]^{w_0}$ lies in
  $\mathcal O_{\mathrm{loc}}(\mathcal E_w)$ with bound (ii).
  Its universal property is precisely the one proved in
  Proposition~\ref{prop:strict-ml-habitat-exactness}: it is the unique
  compatible strict ML element with these finite-window projections.

  The endpoint obstruction class
  $[C_n[\Gamma]]\in H^1(\mathcal O_{\mathrm{loc}}(\mathcal E_w),
  Q+\{I_0,-\})$ is represented by the inverse limit of the compatible
  windowwise classes when that compatibility holds.  The theorem
  identifies its analytic support and coefficient topology.  Its
  vanishing is an additional algebraic/BV condition on the mixed
  brane-defect interaction.
\end{proof}

\begin{rmk}[Why the Hörmander commutation works]\label{rmk:hormander-commutation}
  The deciding technical point is whether the Mittag-Leffler limit
  in the coefficient direction commutes with Hörmander's
  wave-front-set analysis on the bulk. Three facts make this work:

  (a) The propagator factorization
  \eqref{eq:wt-propagator-factor} is exact (not asymptotic), so
  the bulk wave-front set is computed once and tensored.

  (b) The coefficient factor $C_{\mathfrak h,w}$ is not a
  distribution on any spacetime: it is an algebraic tensor in the
  weighted Tate category. Wave-front analysis on a smooth/algebraic
  tensor is trivial ($\mathrm{WF}=\emptyset$).

  (c) The Hadamard parametrix coefficients $a_k(x,y)$ are universal
  in the bulk (depend only on the bulk Riemannian/Hermitian data of
  $\R^2\times\C^2$); they do not see the coefficient direction.

  Consequently, the wave-front bound of
  Lemma~\ref{lem:windowwise-parametrix}~(3) is exact and
  $w_0$-uniform; passing to the Mittag-Leffler limit preserves the
  local-functional class graph-by-graph at every order $\hbar^n$ once
  the windowwise counterterm system is compatible.
\end{rmk}

\begin{cor}[Hadamard coefficient criterion for the weighted QME lift]\label{cor:wt-cotangent-lift-hadamard-criterion}
  In the graphwise Mittag-Leffler sector of
  Theorem~\ref{thm:hadamard-mittag-leffler},
  Proposition~\ref{prop:wt-qme-lift-obstruction-vanishing} is equivalent to the
  vanishing of the mixed brane-defect QME obstruction class.  The
  all-order quantum source-to-boundary CE/PV comparison
  \[
  \begin{aligned}
    \Phi_\hbar^{w}\colon
    &C^\bullet_{\mathrm{CE},\mathrm{cont}}\bigl(
      \mathfrak h_{w,\hbar}\ltimes
      \mathfrak h^{\vee,w}_{\hbar,\mathrm{cont}}[1]\bigr)[[\hbar]]\\
    &\longrightarrow
      Z^{P_{0,\hbar}}_{\mathrm{fact}}\bigl(
      \mathrm{Obs}^\hbar_{\mathrm{open}}\bigr)\big|_{\mathfrak g_w}
  \end{aligned}
  \]
  extends to the weighted cotangent direction at all orders in $\hbar$
  as a continuous $\C[[\hbar]]$-linear cochain map precisely when the
  obstruction class in
  \(H^1(\mathcal O_{\mathrm{loc}}(\mathcal E_w),Q+\{I_0,-\})\)
  vanishes.  Equivalently, this vanishing statement is the datum required
  by Corollary~\ref{cor:wt-descendant-qme-vanishing}.
\end{cor}

\begin{proof}
  Proposition~\ref{prop:wt-qme-lift-obstruction-vanishing} assumes the
  Costello-locality/QME compatibility identity of
  Problem~\ref{prob:weighted-rg-locality}. By
  Theorem~\ref{thm:hadamard-mittag-leffler}, the Hadamard
  wave-front-set and Mittag-Leffler coefficient-topology pieces of
  that identity are controlled.  The proposition therefore asks for
  exactly the vanishing of the mixed brane-defect QME obstruction class.
\end{proof}

\begin{thm}[Weighted RG locality obstruction criterion]\label{thm:weighted-rg-locality-reduced}
  Before adjoining the principal-part current extension, in the graphwise
  finite-window sector of
  Theorem~\ref{thm:hadamard-mittag-leffler}, the analytic part of
  Problem~\ref{prob:weighted-rg-locality} is exactly the following
  finite-window coefficient criterion:
  \[
  \begin{aligned}
    P_{\epsilon,L}^{w}
      &=P_{\epsilon,L}^{\mathrm{base}}\widehat\otimes
        C_{\mathfrak h\oplus\mathfrak h^\vee[1],w},\\
    \mathrm{WF}(C_{\mathfrak h,w}^{w_0})
      &=\emptyset,\\
    \varprojlim\nolimits^1_{w_0}P_{\epsilon,L}^{w,w_0}
      &=0.
  \end{aligned}
  \]
  These conditions hold by Lemma~\ref{lem:windowwise-parametrix}.  After
  they are imposed, the endpoint datum in the weighted QME lift is the
  class
  \[
    \operatorname{Ob}_{w}\in
    H^1(\mathcal O_{\mathrm{loc}}(\mathcal E_w),Q+\{I_0,-\}).
  \]
\end{thm}

\begin{proof}
  The propagator factorization is \eqref{eq:wt-propagator-factor}; its
  finite-window form is the definition of \(P_{\epsilon,L}^{w,w_0}\) in
  Lemma~\ref{lem:windowwise-parametrix}.  The same lemma proves that the
  coefficient Casimir is an algebraic tensor with empty spacetime
  wave-front set and that the finite-window propagators form a strict
  Mittag-Leffler system.  Thus the Hadamard singular support, asymptotic
  coefficients, and locality estimates are exactly the bulk Costello
  data tensored with a compatible coefficient kernel.  The only term of
  Problem~\ref{prob:weighted-rg-locality} not decided by these analytic
  statements is the local-functional QME class displayed above. This does
  not decide the \(\mathcal D'_c(I)\otimes\mathcal P[1]\)
  defect-current compatibility clause of
  Problem~\ref{prob:analytic-graph-realization}.
\end{proof}

\subsection*{Finite-window graph/QME assembly}

\begin{defn}[Finite-window graph/QME package]\label{defn:finite-window-graph-qme-package}
  Fix an interval \(I\), an admissible weight \(w\), a finite Hamiltonian
  window \(M\), and a finite set \(\mathcal G_M\) of connected Costello
  graphs.  A finite-window graph/QME package over \((I,w,M,\mathcal G_M)\)
  consists of the following data.
  \begin{enumerate}[label=(F\arabic*)]
    \item For every \(\Gamma\in\mathcal G_M\) with \(r_\Gamma\) cotangent
    external labels, a tuple
    \(B^\Gamma_\bullet=(B^\Gamma_1,\ldots,B^\Gamma_{r_\Gamma})\) satisfying
    one of the two proved label conditions:
    \[
      B^\Gamma_\bullet\in
      \bigl(\mathcal D^{\mathrm{reg}}_c(I)\bigr)^{r_\Gamma}
      \quad\text{or}\quad
      B^\Gamma_\bullet\in
      \mathcal D'_{c,\mathrm{WF}}(I;\Gamma,M).
    \]
    The second condition is the graphwise relation of
    Definition~\ref{def:app-wavefront-admissible-defect-labels}.  If one
    wants an honest CE source from a linear subspace
    \(\mathcal V\subset\mathcal D'_c(I)\), every finite tuple from
    \(\mathcal V\) must satisfy the corresponding graphwise relation.
    Otherwise the source is only the finite CE span of the specified
    labelled inputs.  Denote the resulting source CE object by
    \(C^\bullet_{\mathrm{lab}}(\mathcal G_M)\).

    \item The finite-scale regulator package
    \[
      P_{\epsilon,L}^{w,M},\qquad \Delta_{\epsilon,L}^{w,M},
      \qquad I_{0,M}^{w},
    \]
    together with graph weights
    \(W^M_{\Gamma,\epsilon,L}(B^\Gamma_\bullet)\), local counterterms
    \(C^M_{\Gamma,w}(\epsilon;B^\Gamma_\bullet)\), and renormalized
    weights
    \[
      W^{R,M}_{\Gamma,L,w}
        =
      \lim_{\epsilon\to0}
      \bigl(
        W^M_{\Gamma,\epsilon,L}
        -
        C^M_{\Gamma,w}(\epsilon;B^\Gamma_\bullet)
      \bigr).
    \]
    These are the regular-density weights of
    Proposition~\ref{prop:app-regular-density-costello-graph-kernel}, or
    the wavefront-admissible weights of
    Proposition~\ref{prop:app-wavefront-admissible-costello-graph-kernel}.

    \item A complete filtered local-functional subcomplex
    \[
      \mathfrak Q^\bullet_{\mathcal G,M}(I)
        =
      \left(
        \mathcal O^{\mathrm{bd}}_{\mathrm{loc},\mathcal G,M}\bigl(
          \mathcal E_w^M|_I;\,
          A_{\partial,\hbar}^{\mathrm{pp},w,M}(I)\bigr),
        d_M
      \right),
      \qquad
      d_M=Q+\{I^w_{0,M},-\}_{\hbar,M},
    \]
    generated, and then completed in the \(\hbar\)-adic and scalar-contact
    filtrations, by the displayed renormalized graph weights, their local
    counterterms, and the \(d_M\)-boundaries and BV brackets forced by the
    finite graph differential on \(\mathcal G_M\).  This is a subcomplex
    only after those boundary and bracket terms have been included.

    \item A filtered chain projection on this subcomplex,
    \[
      \widehat\sigma_{\mathrm{sc},M}\colon
      \mathfrak Q^\bullet_{\mathcal G,M}(I)
      \longrightarrow
      C^\bullet_{\mathrm{Lie}}(\bar A;\C\gamma_{\mathbf 1})[1][[\hbar]],
    \]
    whose associated graded is the scalar-symbol projection.  Equivalently,
    either the graph package lies in the balanced graph/counterterm
    subhabitat of
    Definition~\ref{def:app-balanced-costello-graph-subhabitat}, or the
    scalar-extension obstructions
    \[
      \delta^{(0)}_{\mathrm{Cost},\sigma}(I),
      \qquad
      \mathfrak o^{(r)}_{\mathrm{Cost},\sigma}(I),\quad r>0,
    \]
    of Theorem~\ref{thm:app-costello-graph-extension-obstruction} have been
    killed on the chosen finite-window graph/counterterm subcomplex.

    \item A degree-zero graph kernel
    \[
      \kappa^M_{\mathcal G,w,I}\in
      \operatorname{Hom}_{\mathrm{cont}}\!\left(
        C^\bullet_{\mathrm{lab}}(\mathcal G_M),
        \mathfrak Q^\bullet_{\mathcal G,M}(I)\right),
    \]
    restricted to the chosen label package, whose zero-internal-edge part
    is \(\kappa^{\mathrm{coef},M}_{\hbar,w,I}\) of
    Proposition~\ref{prop:app-bulk-defect-kernel-layer} and whose positive
    graph components are the renormalized weights
    \(W^{R,M}_{\Gamma,L,w}\).
  \end{enumerate}
  No summation over graph topologies is part of the datum unless a separate
  graph-summability hypothesis is supplied.
\end{defn}

\begin{thm}[Finite-window graph/QME assembly]\label{thm:finite-window-graph-qme-assembly}
  Let
  \((I,w,M,\mathcal G_M)\) carry a finite-window graph/QME package in the
  sense of Definition~\ref{defn:finite-window-graph-qme-package}.  Put
  \[
    S^\bullet=
    C^\bullet_{\mathrm{Lie}}(\bar A;\C\gamma_{\mathbf 1})[1][[\hbar]],
    \qquad
    \mathcal Q^\bullet_{\mathcal G,M}(I)
      =
    \ker(\widehat\sigma_{\mathrm{sc},M})
      \subset \mathfrak Q^\bullet_{\mathcal G,M}(I).
  \]
  Let
  \[
    \operatorname{Ob}^{\mathrm{graph}}_{\mathcal G,M}
      =
    \operatorname{ev}_{\mathcal G}
    \operatorname{Curv}_{\mathbf K}(\kappa^M_{\mathcal G,w,I})
      \in \mathfrak Q^1_{\mathcal G,M}(I)
  \]
  be the graph curvature evaluated on the chosen finite labelled CE
  cycles; before evaluation it is the corresponding Hom-valued curvature.
  The source-differential term is included in
  \(\operatorname{Curv}_{\mathbf K}\).  Here
  \(\operatorname{Curv}_{\mathbf K}=d_{\mathbf K}
  +\frac12[-,-]_{\mathbf K}\) as in
  Definition~\ref{def:app-bulk-defect-convolution-habitat}.  Suppose
  counterterms have been chosen through order \(\hbar^{n-1}\), and write
  \(C_{<n,M}\in\mathfrak Q^0_{\mathcal G,M}(I)\) for their sum.  The
  order-\(n\) residual is
  \[
    \mathfrak o^{\mathrm{fw}}_{n,M}
      =
    [\hbar^n]\left(
      \operatorname{Ob}^{\mathrm{graph}}_{\mathcal G,M}
      +d_M C_{<n,M}
      +\frac12\{C_{<n,M},C_{<n,M}\}_{\hbar,M}
    \right).
  \]
  Then the following statements hold.
  \begin{enumerate}[label=(\arabic*)]
    \item The analytic graph weights in the package are theorem-grade
    finite-window local functionals.  For regular density labels this is
    Proposition~\ref{prop:app-regular-density-costello-graph-kernel}; for
    wavefront-admissible labels this is
    Proposition~\ref{prop:app-wavefront-admissible-costello-graph-kernel}.
    No statement for arbitrary \(\mathcal D'_c(I)\)-labels follows.

    \item The first obstruction is scalar.  If
    \[
      \mathfrak s_{n,M}:=
      [\hbar^n]\,
      \widehat\sigma_{\mathrm{sc},M}\!\left(
        \operatorname{Ob}^{\mathrm{graph}}_{\mathcal G,M}
        +d_M C_{<n,M}
        +\frac12\{C_{<n,M},C_{<n,M}\}_{\hbar,M}
      \right)
    \]
    is nonzero in \(S^1\), the order-\(n\) graph/QME equation cannot be
    solved inside the non-scalar kernel.  In the balanced graph/counterterm
    subhabitat this scalar term is zero by
    Proposition~\ref{prop:app-balanced-costello-graph-preservation}.

    \item If the scalar term vanishes through order \(n\) and the
    lower-order QME equations hold, then
    \[
      \mathfrak o^{\mathrm{fw}}_{n,M}
        \in Z^1\bigl(\mathcal Q^\bullet_{\mathcal G,M}(I)\bigr).
    \]
    An order-\(n\) finite-window counterterm
    \(C_{n,M}\in\mathcal Q^0_{\mathcal G,M}(I)\) satisfying
    \[
      d_M C_{n,M}=-\mathfrak o^{\mathrm{fw}}_{n,M}
    \]
    exists if and only if
    \[
      [\mathfrak o^{\mathrm{fw}}_{n,M}]
        =0
        \quad\text{in}\quad
      H^1\bigl(\mathcal Q^\bullet_{\mathcal G,M}(I),d_M\bigr).
    \]
    If this class is nonzero, it is the precise finite-window non-scalar
    obstruction to the graph/QME realization at order \(\hbar^n\).
  \end{enumerate}
\end{thm}

\begin{proof}
  The first assertion is exactly the analytic input supplied by the
  regular-density and wavefront-admissible kernel propositions.  The
  finite-window hypothesis makes every coefficient factor finite rank, so
  Theorem~\ref{thm:hadamard-mittag-leffler} supplies the Hadamard
  counterterm control for each fixed graph.  The label restriction is
  essential: outside
  \(\mathcal D'_{c,\mathrm{WF}}(I;\Gamma,M)\), the product of the graph
  Hadamard coefficient with the pushed-forward current label is not a
  canonical distribution.

  The scalar statement is formal.  A filtered chain projection
  \(\widehat\sigma_{\mathrm{sc},M}\) is part of the package.  If the
  displayed scalar projection of the residual is nonzero, then the
  residual does not lie in
  \(\ker\widehat\sigma_{\mathrm{sc},M}\), so no counterterm inside that
  kernel can solve the equation.  In the balanced graph/counterterm
  subhabitat the projection is the zero chain map, hence the scalar branch
  vanishes before the non-scalar class is formed.

  Finally, the completed convolution Hom is a filtered dg Lie algebra.
  The Bianchi identity
  \[
    d_{\mathbf K}\operatorname{Curv}_{\mathbf K}(\kappa)
      +[\kappa,\operatorname{Curv}_{\mathbf K}(\kappa)]_{\mathbf K}=0
  \]
  and the lower-order QME equations imply that the first unsolved
  \(\hbar^n\)-coefficient is closed in the convolution differential.
  After evaluation on labelled CE cycles this is \(d_M\)-closed.
  Since the scalar projection
  is a chain map and its order-\(n\) value is zero, this coefficient lies
  in the kernel complex \(\mathcal Q^\bullet_{\mathcal G,M}(I)\).  Adding a
  degree-zero counterterm changes the first unsolved coefficient by the
  \(d_M\)-boundary of that counterterm; quadratic counterterm terms are
  already included in the displayed residual.  Therefore the finite-window
  equation is solvable exactly when the displayed degree-one cohomology
  class vanishes.
\end{proof}

\begin{prop}[Finite-window-to-limit obstruction]\label{prop:finite-window-qme-limit-condition}
  Suppose that finite-window graph/QME packages
  \(\{(I,w,M,\mathcal G_M)\}_{M\geq1}\) are given and satisfy the
  following additional hypotheses.
  \begin{enumerate}[label=(L\arabic*)]
    \item The graph package is compatible under truncation:
    \(\pi_{M,N}\mathfrak Q^\bullet_{\mathcal G,M}(I)\subset
    \mathfrak Q^\bullet_{\mathcal G,N}(I)\), the kernels, counterterms,
    and scalar projections commute with \(\pi_{M,N}\), and the same
    compatibility holds under admissible weight transport \(R^M_{w,w'}\).
    \item At the fixed \(\hbar\)-order under consideration the graph sum is
    finite in each window, or a separate Costello graph-summability datum
    has been supplied.  The finite-window theorem does not create this
    datum.
    \item The residuals
    \(\mathfrak o^{\mathrm{fw}}_{n,M}\) of
    Theorem~\ref{thm:finite-window-graph-qme-assembly} are compatible:
    \[
      \pi_{M,N}\mathfrak o^{\mathrm{fw}}_{n,M}
        =
      \mathfrak o^{\mathrm{fw}}_{n,N},
      \qquad M\geq N.
    \]
  \end{enumerate}
  Put
  \[
    \mathfrak O^{\mathrm{fw}}_n
      =
    \left(
      \bigl([\mathfrak o^{\mathrm{fw}}_{n,M}]\bigr)_M,\,
      \lambda^{\mathrm{fw}}_n
    \right),
  \]
  where
  \[
    \bigl([\mathfrak o^{\mathrm{fw}}_{n,M}]\bigr)_M
      \in
    \varprojlim_M
    H^1\bigl(\mathcal Q^\bullet_{\mathcal G,M}(I)\bigr),
  \]
  and, after choosing windowwise primitives
  \(c_{n,M}\) with
  \(d_M c_{n,M}=-\mathfrak o^{\mathrm{fw}}_{n,M}\),
  \[
    \lambda^{\mathrm{fw}}_n\in
    \varprojlim\nolimits^1_M
    H^0\bigl(\mathcal Q^\bullet_{\mathcal G,M}(I)\bigr)
  \]
  is represented by the Roos cocycle
  \[
    \bigl[\pi_{M,N}c_{n,M}-c_{n,N}\bigr]\in
    H^0\bigl(\mathcal Q^\bullet_{\mathcal G,N}(I)\bigr).
  \]
  Then compatible order-\(n\) counterterms
  \(C_{n,M}\) with
  \(\pi_{M,N}C_{n,M}=C_{n,N}\) exist if and only if
  \(\mathfrak O^{\mathrm{fw}}_n=0\).  If the tower
  \(\{H^0(\mathcal Q^\bullet_{\mathcal G,M}(I))\}_M\) is
  Mittag--Leffler, then \(\lambda^{\mathrm{fw}}_n=0\), and the inverse-limit
  obstruction is only the compatible \(H^1\)-class above.

  Without (L1)--(L3), no inverse-limit graph/QME theorem is asserted.  If
  (L2) fails, the graph sum itself is not an element of the tower, so the
  obstruction class is not formed.
\end{prop}

\begin{proof}
  Windowwise, Theorem~\ref{thm:finite-window-graph-qme-assembly} says that
  the only non-scalar obstruction is the class of
  \(\mathfrak o^{\mathrm{fw}}_{n,M}\) in \(H^1\).  Necessity of the two
  displayed inverse-limit conditions is immediate from any compatible
  primitive \(C_{n,M}\).  Conversely, vanishing of the compatible
  \(H^1\)-class gives primitives \(c_{n,M}\) window by window.  The
  differences \(\pi_{M,N}c_{n,M}-c_{n,N}\) are closed degree-zero
  elements, and their cohomology classes are the standard Roos
  \(1\)-cocycle computing
  \(\varprojlim\nolimits^1_M
  H^0(\mathcal Q^\bullet_{\mathcal G,M}(I))\).  The
  vanishing of \(\lambda^{\mathrm{fw}}_n\) is exactly the statement that
  the \(c_{n,M}\)'s can be adjusted by closed degree-zero terms and
  boundaries to become compatible under all truncation maps.  The
  Mittag--Leffler sufficient condition is the standard vanishing of
  \(\varprojlim^1\) for a Mittag--Leffler tower.  The weight-transport
  compatibility follows from the same diagonal transport squares used in
  Theorem~\ref{thm:wt-regulator-independence-admissible}.
\end{proof}

\begin{defn}[Finite-window graph array]\label{defn:finite-window-graph-array}
  Fix a finite-window graph/QME package
  \((I,w,M,\mathcal G_M)\), a codimension-two closed marked Costello list
  \(\mathcal G^{\mathrm{mk}}_M\), and a labelled CE cycle \(\zeta_M\).
  A finite-window graph array at order \(n\) is a finite set
  \[
    \mathsf A^{\mathrm{fw}}_{n,M}
      =
    \{\alpha\}
  \]
  together with the following row data.

  \begin{enumerate}[label=(A\arabic*)]
    \item \emph{Graph type and order.}  A row has a type
    \[
      \tau_\alpha\in
      \{\mathrm{dK},\mathrm{CE},\mathrm{bracket},
        \mathrm{lower}\text{-}d,\mathrm{lower}\text{-}\mathrm{bracket}\}.
    \]
    For \(\mathrm{dK}\), \(\mathrm{CE}\), and \(\mathrm{bracket}\) rows it
    specifies, respectively, a marked graph \(\Gamma_\alpha\), a CE-source
    face of such a graph, or an ordered pair
    \((\Gamma_{\alpha,1},\Gamma_{\alpha,2})\), all in the finite marked list.
    The order is
    \[
    \begin{array}{c|c}
      \tau_\alpha & |\alpha|_\hbar \\ \hline
      \mathrm{dK},\mathrm{CE} & |\Gamma_\alpha|_\hbar\\
      \mathrm{bracket} &
        |\Gamma_{\alpha,1}|_\hbar+|\Gamma_{\alpha,2}|_\hbar\\
      \mathrm{lower}\text{-}d,\mathrm{lower}\text{-}\mathrm{bracket} & n ,
    \end{array}
    \]
    and only rows with \(|\alpha|_\hbar=n\) enter
    \(\mathsf A^{\mathrm{fw}}_{n,M}\).

    \item \emph{Edge labels.}  Each internal edge of the graph carries one of
    the finite-window edge kernels
    \[
      P_{\epsilon,L}^{w,M},\qquad
      \Delta_{\epsilon,L}^{w,M},\qquad
      E_{\mathrm{Weyl}},
    \]
    where \(E_{\mathrm{Weyl}}\) is the boundary Weyl edge whose scalar
    two-Hamiltonian part is
    \((f,g)\mapsto\omega(f,g)\gamma_{\mathbf 1}\).  The two half-edge
    coefficient labels lie in
    \[
      \mathfrak h_{w,\leq M}\cup D^w_{\leq M},
      \qquad
      D^w_{\leq M}=\prod_{d\leq M}w(d)^{-1}H_d^\vee,
    \]
    and the edge record includes the source/target orientation used by the
    BV bracket sign.

    \item \emph{Vertex labels.}  Each vertex is labelled by one of the
    finite package labels
    \[
      u_{a(t)dt\otimes f},\quad c_{B,\rho},\quad I^w_{0,M},\quad
      W^{R,M}_{\Gamma,L,w}(B^\Gamma_\bullet),\quad
      C^M_{\Gamma,w}(\epsilon;B^\Gamma_\bullet),\quad
      \gamma_{\mathbf 1},
    \]
    with \(f,\rho\) in degree \(\leq M\) and
    \(B^\Gamma_\bullet\) either regular-density or graphwise
    wavefront-admissible as in
    Definition~\ref{defn:finite-window-graph-qme-package}.

    \item \emph{Markers and incidence coefficient.}  The row contains a
    marker tensor
    \[
      \mathfrak m_\alpha=(\mathfrak m_{\alpha,L})_L,\qquad
      \mathfrak m_{\alpha,L}\in\mathcal M^{\mathrm{full}}_{m_L}(\C^{N|N})
    \]
    or the explicitly chosen smaller marker habitat.  If the row comes from
    an elementary boundary operation
    \(\beta_i\in\mathsf B_M(\Gamma_\alpha)\), its incidence coefficient is
    \[
      \iota_{\alpha,M}=(-1)^{i-1}\varepsilon_{\beta_i}.
    \]
    For CE-source and convolution-bracket rows, \(\iota_{\alpha,M}\) is the
    corresponding CE or BV Koszul sign, including the factor \(1/2\) in the
    bracket row.

    \item \emph{Weight and filtrations.}  The row records
    \[
      d_{\max}(\alpha)=
      \max\{d:\ H_d\text{ or }H_d^\vee\text{ occurs in the row}\}\leq M,
    \]
    its scalar-contact filtration
    \(p_\alpha=\#\{\text{scalar-contact open loops}\}\), and, for a boundary
    operation \(\beta_i\),
    \[
      r_\alpha=p_\alpha-
      p(\partial_{\beta_i}\Gamma_\alpha,\mu_{\beta_i}(\mathfrak m_\alpha)).
    \]

    \item \emph{Curvature value and scalar gate.}  The row evaluates to a
    local functional
    \[
      R_{\alpha,M}^{\mathrm{row}}\in\mathfrak Q^1_{\mathcal G,M}(I).
    \]
    For the graph-curvature part these are
    \[
    \begin{array}{c|c}
      \tau_\alpha & R_{\alpha,M}^{\mathrm{row}} \\ \hline
      \mathrm{dK} & \iota_{\alpha,M}\,d_MK^M_{\Gamma_\alpha}\\
      \mathrm{CE} & \iota_{\alpha,M}\,K^M_{\Gamma_\alpha}\\
      \mathrm{bracket} &
        \iota_{\alpha,M}\{K^M_{\Gamma_{\alpha,1}},
        K^M_{\Gamma_{\alpha,2}}\}_{\hbar,M}.
    \end{array}
    \]
    The two lower-counterterm row types are the corresponding homogeneous
    summands of
    \([\hbar^n]d_MC_{<n,M}\) and
    \(\frac12[\hbar^n]\{C_{<n,M},C_{<n,M}\}_{\hbar,M}\).
    The scalar-gate entry is
    \[
      S_{\alpha,M}=
      \widehat\sigma_{\mathrm{sc},M}(R_{\alpha,M}^{\mathrm{row}}).
    \]
    If \(\alpha\) is a first filtration-lowering two-Hamiltonian scalar
    contact row, this specializes to
    \[
      S_{\alpha,M}=
      \iota_{\alpha,M}\,
      \operatorname{Str}_{\mathrm{cyc}}
        \bigl(\mu_{\beta_i}(\mathfrak m_\alpha)\bigr)\,
      \omega(f_{\partial_{\beta_i}\Gamma_\alpha},
             g_{\partial_{\beta_i}\Gamma_\alpha})
      \gamma_{\mathbf 1}.
    \]
    The non-scalar reduced-curvature contribution is not an independent
    entrywise projection.  It is the scalar-zero sum
    \[
      R^{\mathrm{ns}}_{n,M}
        =
      \sum_{\alpha\in\mathsf A^{\mathrm{fw}}_{n,M}}
        R_{\alpha,M}^{\mathrm{row}}
      \quad\text{provided}\quad
      \sum_{\alpha\in\mathsf A^{\mathrm{fw}}_{n,M}}S_{\alpha,M}=0.
    \]
    If an entrywise non-scalar summand is wanted, the array must also supply a
    chain splitting of \(\widehat\sigma_{\mathrm{sc},M}\); without such a
    splitting only the displayed scalar-zero sum is canonical.

    \item \emph{Truncation in \(M\).}  For \(M\geq N\) the array contains
    truncation coefficients \(t^{M,N}_{\alpha\alpha'}\in\C\) with
    \(\alpha'\in\mathsf A^{\mathrm{fw}}_{n,N}\), defined by
    \[
      \pi_{M,N}R^{\mathrm{row}}_{\alpha,M}
        =
      \sum_{\alpha'}t^{M,N}_{\alpha\alpha'}R^{\mathrm{row}}_{\alpha',N},
      \qquad
      \pi_{M,N}S_{\alpha,M}
        =
      \sum_{\alpha'}t^{M,N}_{\alpha\alpha'}S_{\alpha',N}.
    \]
    A row whose coefficient labels have no degree-\(\leq N\) component has
    all \(t^{M,N}_{\alpha\alpha'}=0\).  If this equation cannot be filled for
    a row, the lower-window theorem is missing exactly that truncation
    entry.

    \item \emph{Order-two boundary compositions.}  For an order-\(2\)
    closure check the array also records the ordered pairs of distinct
    first scalar-contact boundary operations
    \[
      (\beta_i,\beta_j),\qquad
      r(\beta_i)=r(\beta_j)=1,\qquad
      |\beta_i|_\hbar=|\beta_j|_\hbar=1,
    \]
    whose two composites occur in the codimension-two closure table of
    Definition~\ref{def:app-finite-window-marked-costello-list}.  If
    \(\operatorname{pos}_i(j)\) denotes the position of the induced face
    \(\beta_j|\beta_i\) after deleting \(\beta_i\), the composite
    incidence coefficient is
    \[
      \iota_{\beta_j|\beta_i,\beta_i}
        =
      {(-1)}^{i-1}{(-1)}^{\operatorname{pos}_i(j)-1}
      \varepsilon_{\beta_i}\varepsilon_{\beta_j|\beta_i}.
    \]
    The corresponding two-contact scalar gate is the row
    \[
      S^{(2)}_{\beta_j|\beta_i,M}
        =
      \iota_{\beta_j|\beta_i,\beta_i}
      \left(
        \prod_{L\in\mathcal L_{\mathrm{sc}}
        (\partial_{\beta_j|\beta_i}\partial_{\beta_i}\Gamma)}
        \operatorname{Str}_{\mathrm{cyc}}(\mathfrak m'_L)
      \right)
      \omega(f_{\beta_i},g_{\beta_i})
      \omega(f_{\beta_j|\beta_i},g_{\beta_j|\beta_i})
      \gamma_{\mathbf 1}^{(2)}.
    \]
    Here \(\mathfrak m'_L\) is the transported marker on the output scalar
    loop \(L\), and \(\gamma_{\mathbf 1}^{(2)}\) denotes the ordered
    two-contact scalar symbol.  This notation records the associated
    graded scalar-contact symbol; it is not a separate assertion that a
    product \(\gamma_{\mathbf 1}^2\) survives in the target algebra.
    If the equality of output graphs, coefficients, and marker transports
    in~\eqref{eq:app-codim-two-marked-closure} is not supplied, the missing
    order-\(2\) datum is exactly that composite face.
  \end{enumerate}

  The array is \emph{valid} if the row sum reproduces the order-\(n\)
  residual:
  \[
    R_{n,M}
      =
    \sum_{\alpha\in\mathsf A^{\mathrm{fw}}_{n,M}}
      R_{\alpha,M}^{\mathrm{row}},
    \qquad
    \mathfrak s_{n,M}
      =
    \sum_{\alpha\in\mathsf A^{\mathrm{fw}}_{n,M}}S_{\alpha,M},
  \]
  and if the truncation coefficients satisfy the displayed equations for
  every \(M\geq N\).
\end{defn}

\begin{prop}[First finite-window array rows]\label{prop:first-finite-window-array-rows}
  The universal rows that are presently computed are the reduced current
  interface row in order \(0\) and the two-Hamiltonian scalar-contact row in
  order \(1\).

  On the length-one reduced current interface,
  Theorem~\ref{thm:app-factorization-universal-current-interface} gives
  strict CE/PV compatibility.  Thus the order-\(0\) row has
  \[
    R^{\mathrm{row}}_{0,M}=0,\qquad S_{0,M}=0.
  \]

  For the ordered Hamiltonian inputs \((a,z_1)\), \((b,z_2)\), the
  one-cross boundary row has graph type \(\mathrm{dK}\), one boundary Weyl
  edge \(E_{\mathrm{Weyl}}\), two external Hamiltonian vertices, one
  scalar-contact output vertex \(\gamma_{\mathbf 1}\), and
  \[
    \omega(z_1,z_2)=1,\qquad |\alpha_{\mathrm{sc}}|_\hbar=1.
  \]
  Its scalar gate is
  \[
    S_{\alpha_{\mathrm{sc}},M}
      =
    \operatorname{Str}_{\mathrm{cyc}}
      \bigl(\mu_{\beta_{\mathrm{sc}}}(\mathfrak m)\bigr)
    \gamma_{\mathbf 1},
  \]
  after extracting the \(\hbar^1\)-coefficient and using the orientation
  convention in which the ordered pair \((z_1,z_2)\) has positive sign.
  Therefore:
  \[
  \begin{array}{c|c|c}
    \text{branch} & S_{\alpha_{\mathrm{sc}},M} &
    R^{\mathrm{ns}}_{1,M}\text{ from this row}\\ \hline
    \lie{gl}_N\text{ ordinary} & N\gamma_{\mathbf 1} & \text{not formed}\\
    \text{even block marker }\lambda_0P_{\bar0}+\lambda_1P_{\bar1}
      & N(\lambda_0-\lambda_1)\gamma_{\mathbf 1} & \text{not formed unless }
        \lambda_0=\lambda_1\\
    \lie{gl}(N|N)\text{ full-equivariant}
      & 0 & 0
  \end{array}
  \]
  Consequently, for the minimal full-equivariant closed array containing only
  the strict current-interface rows and this scalar-contact row, the first
  order residual is zero and
  \[
    C_{1,M}=0
  \]
  is the compatible order-one non-scalar counterterm.  For any larger graph
  package, the first non-scalar residual is
  \[
    R^{\mathrm{ns}}_{1,M}
      =
    \sum_{\alpha\in
      \mathsf A^{\mathrm{fw}}_{1,M}\setminus\{\alpha_{\mathrm{sc}}\}}
      R^{\mathrm{row}}_{\alpha,M}
  \]
  after the scalar gate has vanished.  Its primitive cannot be decided
  without the remaining order-one rows of
  Definition~\ref{defn:finite-window-graph-array}.
\end{prop}

\begin{proof}
  The order-\(0\) assertion is exactly the strict current-source/target
  identity of Theorem~\ref{thm:app-factorization-universal-current-interface}
  on the length-one generators: the CE differential and target \(P_0\)
  bracket use the same semidirect product constants.

  The order-\(1\) scalar-contact row is the computation of
  Proposition~\ref{prop:app-first-nonidentity-loop-scalar-symbol}, written in
  the row notation of Definition~\ref{defn:finite-window-graph-array}.  The
  monomial value is \(\omega(z_1,z_2)=1\).  In the ordinary branch the open
  index loop contributes \(N\).  In the even block-diagonal marker branch
  the cyclic supertrace is \(N(\lambda_0-\lambda_1)\).  In the
  full-equivariant branch, Lemma~\ref{lem:app-full-equivariant-cyclic-supertrace}
  kills the cyclic supertrace of every transported marker tensor in
  \(\C^{N|N}\); the empty marker word is killed by the same
  \(\operatorname{sdim}\C^{N|N}=0\) calculation.  Thus the row contributes
  zero as a local functional in the full-equivariant branch, not merely zero
  after cohomology.

  If no further order-one rows are present, the order-one residual is zero,
  so the zero counterterm is a primitive and is compatible under all
  truncation maps.  If further order-one rows are present, the displayed sum
  is the definition of the remaining scalar-zero residual.  A primitive for
  it is a theorem exactly when the missing row values, signs, scalar gates,
  and truncation coefficients are supplied and the resulting \(H^1\)-class
  vanishes.
\end{proof}

\begin{prop}[Universal scalar-contact bracket rows]\label{prop:universal-scalar-contact-bracket-rows}
  Work in the local full-equivariant marked branch with \(V=\C^{N|N}\), and
  let \(\alpha_{\mathrm{sc}}\) be the one-cross two-Hamiltonian
  scalar-contact row of
  Proposition~\ref{prop:first-finite-window-array-rows}.  Write
  \(K^M_{\alpha_{\mathrm{sc}}}\) for the local functional carried by this
  row.  Then
  \[
    K^M_{\alpha_{\mathrm{sc}}}=0
  \]
  in the finite-window local-functional complex.

  Consequently every order-\(2\) convolution-bracket row containing
  \(\alpha_{\mathrm{sc}}\) has the following entries.  For any order-\(1\)
  component \(\beta\) for which the bracket row is present,
  \[
  \begin{array}{c|c}
    \text{row} & R^{\mathrm{row}}_{b,M} \\ \hline
    b(\alpha_{\mathrm{sc}},\beta) &
      \frac12\varepsilon_{\alpha_{\mathrm{sc}},\beta}
      \{K^M_{\alpha_{\mathrm{sc}}},K^M_{\beta}\}_{\hbar,M}=0\\
    b(\beta,\alpha_{\mathrm{sc}}) &
      \frac12\varepsilon_{\beta,\alpha_{\mathrm{sc}}}
      \{K^M_{\beta},K^M_{\alpha_{\mathrm{sc}}}\}_{\hbar,M}=0\\
    b(\alpha_{\mathrm{sc}},\alpha_{\mathrm{sc}}) &
      \frac12\varepsilon_{\alpha_{\mathrm{sc}},\alpha_{\mathrm{sc}}}
      \{K^M_{\alpha_{\mathrm{sc}}},
        K^M_{\alpha_{\mathrm{sc}}}\}_{\hbar,M}=0
  \end{array}
  \]
  Their scalar gates are
  \[
    S_{b,M}=0.
  \]
  For every \(M\geq N\), the truncation row may be chosen as
  \[
    t^{M,N}_{b b'}=0
    \qquad
    \text{for every lower-window bracket row } b',
  \]
  since \(\pi_{M,N}R^{\mathrm{row}}_{b,M}=0\) and
  \(\pi_{M,N}S_{b,M}=0\).  These rows therefore contribute no order-\(2\)
  obstruction and require no counterterm.

  This proposition closes only the bracket family containing
  \(\alpha_{\mathrm{sc}}\).  Bracket rows \(b(\beta_1,\beta_2)\) with
  \(\beta_1,\beta_2\neq\alpha_{\mathrm{sc}}\) remain part of the chosen
  graph package and must be supplied row by row.
\end{prop}

\begin{proof}
  Proposition~\ref{prop:first-finite-window-array-rows} proves the stronger
  full-equivariant fact that the scalar-contact row is zero as a local
  functional, because every transported scalar-loop marker has zero cyclic
  supertrace on \(\C^{N|N}\).  The BV/convolution bracket is bilinear and
  continuous on the finite-window completed local-functional complex, hence
  any bracket with \(K^M_{\alpha_{\mathrm{sc}}}=0\) is zero before passing to
  cohomology.  Applying the chain projection
  \(\widehat\sigma_{\mathrm{sc},M}\) to zero gives the displayed scalar
  gates.  Truncation sends zero to zero, so the zero truncation matrix row is
  compatible with the defining equations of
  Definition~\ref{defn:finite-window-graph-array}.
\end{proof}

\begin{prop}[Genuine order-two graph-weight rows]\label{prop:genuine-order-two-graph-weight-rows}
  Work in the local
  \(\R^2_{\mathrm{top}}\times\C^2_{\mathrm{hol}}\) full-equivariant marked
  branch with \(V=\C^{N|N}\).  Let
  \(\mathcal G^{\mathrm{gen}}_{2,M}\) be the part of a finite marked
  Costello list consisting of genuine seed graph weights of
  \(\hbar\)-order \(2\): these are not the codimension-two composites of
  two first scalar-contact boundary faces and not the bracket rows
  containing \(\alpha_{\mathrm{sc}}\) closed in
  Proposition~\ref{prop:universal-scalar-contact-bracket-rows}.  For
  \(\Gamma\in\mathcal G^{\mathrm{gen}}_{2,M}\), write
  \[
    K^M_\Gamma=
    W^{R,M}_{\Gamma,L,w}(B^\Gamma_\bullet)
  \]
  for the normalized renormalized finite-window graph weight, and let
  \(\varepsilon_\Gamma\) be the coefficient with which this component
  occurs in \(d_M\kappa^M_{\mathcal G,w,I}\), including the orientation,
  Koszul, and symmetry factors supplied by the finite graph package.

  The genuine order-two rows have the following schema.
  \begin{enumerate}[label=(G\arabic*)]
    \item The graph-differential row \(d\Gamma\) has
      \[
        \tau_{d\Gamma}=\mathrm{dK},\qquad |d\Gamma|_\hbar=2,
      \]
      inherits its edge labels, vertex labels, marker tensor, and
      coefficient degree from \(\Gamma\), and has
      \[
        R^{\mathrm{row}}_{d\Gamma,M}
          =
        \varepsilon_\Gamma\,d_M K^M_\Gamma,
        \qquad
        S_{d\Gamma,M}
          =
        \varepsilon_\Gamma\,
        \widehat\sigma_{\mathrm{sc},M}(d_M K^M_\Gamma).
      \]

    \item Before evaluation on a strict CE cycle, or when the Hom-valued
      kernel table is retained, each source face is a separate row.  Write
      the relevant part of the source differential as
      \[
        d_{\mathrm{CE}}\zeta_M
          =
        \sum_\nu
        \varepsilon^{\mathrm{CE}}_{\Gamma,\nu}\zeta_{M,\nu},
      \]
      where \(\varepsilon^{\mathrm{CE}}_{\Gamma,\nu}\) includes the CE
      Koszul sign and the graph coefficient attached to the
      \(\Gamma\)-component.  The source row is
      \[
        R^{\mathrm{row}}_{\mathrm{CE}(\Gamma,\nu),M}
          =
        -\varepsilon^{\mathrm{CE}}_{\Gamma,\nu}
        K^M_\Gamma(\zeta_{M,\nu}),
      \]
      and its scalar gate is
      \[
        S_{\mathrm{CE}(\Gamma,\nu),M}
          =
        -\varepsilon^{\mathrm{CE}}_{\Gamma,\nu}
        \widehat\sigma_{\mathrm{sc},M}
        \bigl(K^M_\Gamma(\zeta_{M,\nu})\bigr).
      \]
      If the curvature has already been evaluated on a strict labelled CE
      cycle, this source-face family is empty.
  \end{enumerate}

  If these row values and all their \(d_M\)-faces lie in the
  full-equivariant marked habitat of
  Theorem~\ref{thm:app-full-equivariant-marked-shadow-vanishing}, then
  \[
    S_{d\Gamma,M}=0,
    \qquad
    S_{\mathrm{CE}(\Gamma,\nu),M}=0.
  \]
  This is only a scalar-gate statement.  It does not imply
  \(K^M_\Gamma=0\) or \(d_M K^M_\Gamma=0\).

  After scalar-gate vanishing, the genuine seed-graph contribution to the
  order-\(2\) reduced curvature is
  \[
  \begin{aligned}
    R^{\mathrm{ns,gen}}_{2,M}
      &=
      \sum_{\Gamma\in\mathcal G^{\mathrm{gen}}_{2,M}}
        \varepsilon_\Gamma\,d_M K^M_\Gamma
      \\
      &\quad -
      \sum_{\Gamma,\nu}
        \varepsilon^{\mathrm{CE}}_{\Gamma,\nu}
        K^M_\Gamma(\zeta_{M,\nu})
      \in
      \mathcal K^1_{\mathrm{ns},M}(I).
  \end{aligned}
  \]
  The second sum is absent after strict CE-cycle evaluation.  A genuine
  order-\(2\) counterterm exists precisely when
  \[
    [R^{\mathrm{ns,gen}}_{2,M}]
      =
    0
    \quad\text{in}\quad
    H^1(\mathcal K^\bullet_{\mathrm{ns},M}(I),d_{\mathrm{ns},M}).
  \]
  If this class is nonzero, it is the explicit finite-window obstruction
  carried by the genuine seed graphs.

  Truncation is part of the row data.  For \(M\geq N\) one must supply
  normalized coefficients
  \[
    \pi_{M,N}(\varepsilon_\Gamma K^M_\Gamma)
      =
    \sum_{\Gamma'}
      u^{M,N}_{\Gamma\Gamma'}\,\varepsilon_{\Gamma'}K^N_{\Gamma'},
  \]
  and then the graph-differential truncation entries are
  \[
    t^{M,N}_{d\Gamma,d\Gamma'}=u^{M,N}_{\Gamma\Gamma'}.
  \]
  For source rows one must likewise supply
  \[
    \pi_{M,N}
      \bigl(
        -\varepsilon^{\mathrm{CE}}_{\Gamma,\nu}
        K^M_\Gamma(\zeta_{M,\nu})
      \bigr)
      =
    \sum_{\Gamma',\nu'}
      v^{M,N}_{\Gamma,\nu;\Gamma',\nu'}\,
      \bigl(
        -\varepsilon^{\mathrm{CE}}_{\Gamma',\nu'}
        K^N_{\Gamma'}(\zeta_{N,\nu'})
      \bigr),
  \]
  and set
  \(t^{M,N}_{\mathrm{CE}(\Gamma,\nu),\mathrm{CE}(\Gamma',\nu')}=
  v^{M,N}_{\Gamma,\nu;\Gamma',\nu'}\).  Rows whose coefficient labels have
  no degree-\(\leq N\) component have zero truncation coefficients.  If
  either displayed truncation identity is not supplied, the lower-window
  theorem is missing exactly that row entry.

  In the minimal full-equivariant local package of
  Proposition~\ref{prop:order-two-finite-window-boundary-rows},
  \(\mathcal G^{\mathrm{gen}}_{2,M}=\emptyset\) and the source-face set is
  empty.  Hence
  \[
    R^{\mathrm{ns,gen}}_{2,M}=0,
    \qquad
    C^{\mathrm{gen}}_{2,M}=0.
  \]
  In any larger package, the present manuscript does not yet decide
  vanishing, counterterms, or obstruction.  The exact missing data are:
  the finite seed list \(\mathcal G^{\mathrm{gen}}_{2,M}\); the
  renormalized local weights \(K^M_\Gamma\) with regular-density or
  wavefront-admissible current labels; the full \(d_M K^M_\Gamma\)
  expansion with orientation signs; the CE-source decomposition and signs
  unless strict CE-cycle evaluation has removed them; proof that every row
  lies in the full-equivariant marked habitat, or else its scalar
  projection; and the truncation matrices displayed above.
\end{prop}

\begin{proof}
  This is the order-\(2\) part of the convolution-curvature formula in
  Theorem~\ref{thm:actual-finite-window-nonscalar-decision-datum}, with
  the scalar-contact bracket family removed by
  Proposition~\ref{prop:universal-scalar-contact-bracket-rows}.  The
  differential term contributes \(\varepsilon_\Gamma d_M K^M_\Gamma\).  The
  source term in \(d_{\mathbf K}=d_M-(-1)^{|\cdot|}(\cdot)d_{\mathrm{CE}}\)
  contributes the displayed minus sign for a degree-zero kernel.  After
  evaluation on a strict CE cycle, that source term is zero.

  The scalar-gate assertion follows from
  Theorem~\ref{thm:app-full-equivariant-marked-shadow-vanishing}: on the
  full-equivariant marked habitat the filtered scalar projection is the
  zero chain map.  Since this projection only kills the scalar shadow, it
  gives no vanishing theorem for the non-scalar local functional
  \(K^M_\Gamma\).

  The counterterm criterion is exactly
  Definition~\ref{def:app-nonscalar-kernel-row-complex}: once the scalar
  gate is zero, the row sum is a degree-one element of the non-scalar
  kernel complex, and a degree-zero counterterm exists exactly when this
  element is a \(d_{\mathrm{ns},M}\)-boundary.  The truncation formulas are
  the row identities required in
  Definition~\ref{defn:finite-window-graph-array}, written for normalized
  graph and source rows.  In the minimal package no such genuine seed row
  is present, so the sum is empty.
\end{proof}

\begin{prop}[Order-two finite-window boundary rows]\label{prop:order-two-finite-window-boundary-rows}
  Work in the local
  \(\R^2_{\mathrm{top}}\times\C^2_{\mathrm{hol}}\) full-equivariant marked
  branch with \(V=\C^{N|N}\).  Let the minimal array be generated by the
  strict current-interface rows and the one-cross scalar-contact row of
  Proposition~\ref{prop:first-finite-window-array-rows}.  Assume no genuine
  \(|\Gamma|_\hbar=2\) seed graph, no order-\(2\) CE-source face, and no
  order-\(2\) graph counterterm face have been added to the finite marked
  Costello list, and use \(C_{1,M}=0\).

  Then the only order-\(2\) boundary data forced by codimension-two closure
  are the two ordered composites of two first scalar-contact boundary
  operations.  For
  \(o=\beta_1\wedge\beta_2\) they have incidence signs
  \[
    \iota_{\beta_2|\beta_1,\beta_1}
      =
    \varepsilon_{\beta_1}\varepsilon_{\beta_2|\beta_1},
    \qquad
    \iota_{\beta_1|\beta_2,\beta_2}
      =
    -\varepsilon_{\beta_2}\varepsilon_{\beta_1|\beta_2}.
  \]
  Their scalar gates are
  \[
  \begin{aligned}
    S^{(2)}_{\beta_2|\beta_1,M}
      &=
      \varepsilon_{\beta_1}\varepsilon_{\beta_2|\beta_1}
      \prod_L\operatorname{Str}_{\mathrm{cyc}}(\mathfrak m'_L)\,
      \omega(f_{\beta_1},g_{\beta_1})
      \omega(f_{\beta_2|\beta_1},g_{\beta_2|\beta_1})
      \gamma_{\mathbf 1}^{(2)},\\
    S^{(2)}_{\beta_1|\beta_2,M}
      &=
      -\varepsilon_{\beta_2}\varepsilon_{\beta_1|\beta_2}
      \prod_L\operatorname{Str}_{\mathrm{cyc}}(\mathfrak m''_L)\,
      \omega(f_{\beta_2},g_{\beta_2})
      \omega(f_{\beta_1|\beta_2},g_{\beta_1|\beta_2})
      \gamma_{\mathbf 1}^{(2)}.
  \end{aligned}
  \]
  The codimension-two closure equations identify the two output graphs,
  the two products of coefficients, and the transported marker tensors.
  Therefore the pair sum is zero.  In the full-equivariant branch it is
  stronger: every factor
  \(\operatorname{Str}_{\mathrm{cyc}}(\mathfrak m'_L)\) and
  \(\operatorname{Str}_{\mathrm{cyc}}(\mathfrak m''_L)\) is zero by
  Lemma~\ref{lem:app-full-equivariant-cyclic-supertrace}; hence each
  two-contact scalar row already vanishes as a local scalar row.

  Under these minimal hypotheses the order-\(2\) reduced non-scalar row sum
  is
  \[
    R^{\mathrm{ns}}_{2,M}=0,
    \qquad
    C_{2,M}=0
  \]
  is a compatible order-\(2\) counterterm.

  For any larger finite-window package, the order-\(2\) residual is not
  decided by the present local data.  After scalar-gate vanishing its
  exact row sum is
  \[
  \begin{aligned}
    R^{\mathrm{ns}}_{2,M}
     ={}&
      \sum_{|\Gamma|_\hbar=2}
        \varepsilon_\Gamma\,d_M K^M_\Gamma
      -\sum_{|\Gamma'|_\hbar=2}
        \varepsilon^{\mathrm{CE}}_{\Gamma'}K^M_{\Gamma'}\\
      &\quad
      +\frac12
      \sum_{|\Gamma_1|_\hbar=|\Gamma_2|_\hbar=1}
        \varepsilon_{\Gamma_1,\Gamma_2}
        \{K^M_{\Gamma_1},K^M_{\Gamma_2}\}_{\hbar,M}\\
      &\quad
      +[\hbar^2]\left(
        d_M C_{<2,M}
        +\frac12\{C_{<2,M},C_{<2,M}\}_{\hbar,M}\right).
  \end{aligned}
  \]
  Thus a nonzero order-\(2\) obstruction or its counterterm can be asserted
  only after supplying the genuine two-cross/contact graph weights
  \(K^M_\Gamma\), the CE-source faces, the BV bracket rows between
  order-\(1\) graph components not containing
  \(\alpha_{\mathrm{sc}}\) or another retained zero local-functional row
  such as \(\alpha_{11}\), the
  lower-counterterm rows if a nonminimal \(C_{1,M}\) was chosen, the scalar
  projection of any non-full-equivariant two-contact symbol, and the
  truncation coefficients \(t^{M,N}_{\alpha\alpha'}\).
\end{prop}

\begin{proof}
  The signs are those of
  Lemma~\ref{lem:app-oriented-marked-incidence-signs}.  Starting with
  \(o=\beta_1\wedge\beta_2\), the face \(\beta_1\) followed by
  \(\beta_2|\beta_1\) contributes \(+1\), while the face \(\beta_2\)
  followed by \(\beta_1|\beta_2\) contributes \(-1\).  The
  codimension-two closure equations
  Equation~\eqref{eq:app-codim-two-marked-closure} identifies the output graph,
  coefficient, and marker transport in the two composites, so the two
  boundary-composition rows cancel pairwise.

  The scalar-gate formula is the order-\(2\) instance of the last row datum
  in Definition~\ref{defn:finite-window-graph-array}: one scalar-contact
  factor is contributed by each first-lowering boundary face.  In the
  full-equivariant marked habitat, marker transport preserves the signed
  Brauer span.  Lemma~\ref{lem:app-full-equivariant-cyclic-supertrace}
  gives zero cyclic supertrace for every transported scalar-loop marker,
  including the empty marker word.  Hence the two-contact scalar rows
  vanish before passing to cohomology.

  Under the stated minimal hypotheses there are no genuine order-\(2\)
  graph rows, no CE-source rows, no order-\(1\) non-scalar primitive to
  feed a lower-counterterm row, and no nonzero scalar-contact row left
  after the full-equivariant trace calculation.  The possible bracket row
  between two one-cross scalar-contact components is also zero, because
  each such component has zero local-functional value in the
  full-equivariant branch by
  Proposition~\ref{prop:first-finite-window-array-rows}.  The reduced
  non-scalar order-\(2\) row sum is therefore zero, and the zero
  counterterm is compatible under finite-window truncation.

  Removing any minimal hypothesis reintroduces exactly the displayed row
  types from the curvature formula in
  Theorem~\ref{thm:actual-finite-window-nonscalar-decision-datum}.  The
  bracket rows containing \(\alpha_{\mathrm{sc}}\) are the universal zero
  rows of
  Proposition~\ref{prop:universal-scalar-contact-bracket-rows}; bracket rows
  containing a retained \(\alpha_{11}\)-row vanish by the same bilinearity
  argument and
  Proposition~\ref{prop:app-first-scalar-zero-order-one-row}.  Without the
  remaining graph weights, signs, scalar projections, and truncation
  coefficients, the cochain \(R^{\mathrm{ns}}_{2,M}\) is not defined, so no
  order-\(2\) obstruction class or counterterm can be claimed.
\end{proof}

\begin{prop}[Truncation matrices for the computed finite-window rows]\label{prop:computed-finite-window-truncation-matrices}
  Work with the finite-window graph arrays of
  Definition~\ref{defn:finite-window-graph-array}.  For \(M\geq N\), use the
  following label-retaining convention: if a labelled row has
  \(d_{\max}\leq N\), the same label denotes its lower-window row; if
  \(d_{\max}>N\), the row has no lower-window representative and all entries
  from it are zero.  With this convention the computed rows have the
  following truncation matrices.

  \begin{enumerate}[label=(T\arabic*)]
    \item The one-cross scalar-contact row
    \(\alpha_{\mathrm{sc}}\) of
    Proposition~\ref{prop:first-finite-window-array-rows} has
    \(d_{\max}(\alpha_{\mathrm{sc}})=1\).  For \(M\geq N\geq1\),
    \[
      t^{M,N}_{\alpha_{\mathrm{sc}}\alpha_{\mathrm{sc}}}=1,
      \qquad
      t^{M,N}_{\alpha_{\mathrm{sc}}\alpha'}=0
      \quad(\alpha'\neq\alpha_{\mathrm{sc}}).
    \]
    This applies to the ordinary, even-block, and full-equivariant branches.
    In the full-equivariant branch the row value is already zero, so the
    zero matrix would also satisfy the row equations; the displayed matrix is
    the chosen transport of the retained row label.

    \item The labelled scalar-zero row \(\alpha_{11}\) of
    Proposition~\ref{prop:app-first-scalar-zero-order-one-row}, with ordered
    inputs \((z_1,z_1)\), has
    \[
      R^{\mathrm{row}}_{\alpha_{11},M}=0,\qquad
      S_{\alpha_{11},M}=0,
      \qquad d_{\max}(\alpha_{11})=1.
    \]
    If the row convention retains this labelled zero row, then for
    \(M\geq N\geq1\)
    \[
      t^{M,N}_{\alpha_{11}\alpha_{11}}=1,
      \qquad
      t^{M,N}_{\alpha_{11}\alpha'}=0
      \quad(\alpha'\neq\alpha_{11}).
    \]
    If zero rows are suppressed, the \(\alpha_{11}\)-row and its truncation
    matrix are empty.

    \item Let
    \(\gamma_{21}=(\beta_2|\beta_1,\beta_1)\) and
    \(\gamma_{12}=(\beta_1|\beta_2,\beta_2)\) denote the two ordered
    codimension-two scalar-contact composites in
    Proposition~\ref{prop:order-two-finite-window-boundary-rows}.  When the
    degree-\(\leq N\) window contains the two Hamiltonian labels, the
    composite rows are transported diagonally:
    \[
      t^{M,N}_{\gamma_{21}\gamma_{21}}=1,\qquad
      t^{M,N}_{\gamma_{12}\gamma_{12}}=1,\qquad
      t^{M,N}_{\gamma_{ij}\gamma'}=0
      \quad(\gamma'\neq\gamma_{ij}).
    \]
    The codimension-two pair cancellation is therefore stable under
    truncation.  In the full-equivariant branch the stronger rowwise
    vanishing is also stable.

    \item Let \(b\) be any order-\(2\) bracket row containing a retained
    zero local-functional row
    \[
      \alpha_0\in\{\alpha_{\mathrm{sc}},\alpha_{11}\}.
    \]
    Thus \(b=b(\alpha_0,\beta)\), \(b=b(\beta,\alpha_0)\), or
    \(b=b(\alpha_0,\alpha_0)\), whenever the displayed row belongs to the
    chosen finite marked list.  Then
    \[
      R^{\mathrm{row}}_{b,M}=0,\qquad S_{b,M}=0,
    \]
    and one may take the zero truncation row
    \[
      t^{M,N}_{bb'}=0
      \qquad
      \text{for every lower-window bracket row }b'.
    \]

    \item If genuine order-\(2\) seed rows are absent, as in the minimal
    full-equivariant package, their truncation matrices are empty.  If a
    larger package supplies genuine seeds
    \(\Gamma\in\mathcal G^{\mathrm{gen}}_{2,M}\), then the graph-differential
    rows are transported by the normalized matrix \(u^{M,N}\):
    \[
      \pi_{M,N}(\varepsilon_\Gamma K^M_\Gamma)
        =
      \sum_{\Gamma'}
        u^{M,N}_{\Gamma\Gamma'}\,
        \varepsilon_{\Gamma'}K^N_{\Gamma'},
      \qquad
      t^{M,N}_{d\Gamma,d\Gamma'}=u^{M,N}_{\Gamma\Gamma'}.
    \]
    Source-face rows, when present, are transported by the normalized matrix
    \(v^{M,N}\):
    \[
      \pi_{M,N}
      \bigl(-\varepsilon^{\mathrm{CE}}_{\Gamma,\nu}
      K^M_\Gamma(\zeta_{M,\nu})\bigr)
        =
      \sum_{\Gamma',\nu'}
        v^{M,N}_{\Gamma,\nu;\Gamma',\nu'}\,
        \bigl(-\varepsilon^{\mathrm{CE}}_{\Gamma',\nu'}
        K^N_{\Gamma'}(\zeta_{N,\nu'})\bigr),
    \]
    \[
      t^{M,N}_{\mathrm{CE}(\Gamma,\nu),
        \mathrm{CE}(\Gamma',\nu')}
        =
      v^{M,N}_{\Gamma,\nu;\Gamma',\nu'}.
    \]
    Under the projection-closed row-basis hypotheses of
    Proposition~\ref{prop:projection-defined-uvq-truncation-data}, these are
    the coordinate entries of the finite-window projection.  Without those
    hypotheses the lower-window theorem is missing exactly the entries
    \(u^{M,N}_{\Gamma\Gamma'}\) and
    \(v^{M,N}_{\Gamma,\nu;\Gamma',\nu'}\) for those supplied genuine rows.

    \item Bracket rows \(b(\beta_1,\beta_2)\) with
    \(\beta_i\notin\{\alpha_{\mathrm{sc}},\alpha_{11}\}\) are not computed
    by the present local data.  After a lower-window bracket row basis, or
    an equivalent row presentation with chosen coordinates, is fixed, their
    required truncation matrix is the coordinate family \(q^{M,N}\)
    satisfying
    \[
      \pi_{M,N}\!\left(
        \frac12\varepsilon_{\beta_1,\beta_2}
        \{K^M_{\beta_1},K^M_{\beta_2}\}_{\hbar,M}\right)
       =
      \sum_{\beta'_1,\beta'_2}
        q^{M,N}_{\beta_1,\beta_2;\beta'_1,\beta'_2}
        \frac12\varepsilon_{\beta'_1,\beta'_2}
        \{K^N_{\beta'_1},K^N_{\beta'_2}\}_{\hbar,N}.
    \]
    These \(q\)-entries, together with the underlying order-\(1\) row
    values, signs, and lower-window bracket-basis reduction, are still
    missing larger-package data.
  \end{enumerate}

  The nonempty displayed matrices are stable under further truncation:
  for \(M\geq N\geq P\),
  \[
    t^{M,P}=t^{M,N}t^{N,P}
  \]
  on the scalar-contact, scalar-zero, codimension-two, and zero-bracket
  families.  For genuine seed rows and nonzero added bracket rows, stability
  is exactly the cocycle condition
  \[
    u^{M,P}_{\Gamma\Gamma''}
      =
    \sum_{\Gamma'}u^{M,N}_{\Gamma\Gamma'}u^{N,P}_{\Gamma'\Gamma''},
  \]
  \[
    v^{M,P}_{\Gamma,\nu;\Gamma'',\nu''}
      =
    \sum_{\Gamma',\nu'}
      v^{M,N}_{\Gamma,\nu;\Gamma',\nu'}
      v^{N,P}_{\Gamma',\nu';\Gamma'',\nu''},
  \]
  and the analogous equation for \(q\).
\end{prop}

\begin{proof}
  For the one-cross scalar-contact rows the coefficient labels are
  \(z_1,z_2\), and for \(\alpha_{11}\) the labels are \(z_1,z_1\).  These
  all have window degree \(1\).  Therefore \(\pi_{M,N}\) preserves the
  labelled row for \(N\geq1\), and kills it only if the lower window does
  not contain the degree-one labels.  The diagonal matrices in (T1) and
  (T2) are exactly this statement.  The row \(\alpha_{11}\) is zero because
  \(\omega(z_1,z_1)=0\).

  The codimension-two rows use the same two degree-one Hamiltonian labels
  and the same transported marker words before and after projection.  The
  finite-window marked-habitat compatibility preserves the output graph,
  coefficient product, marker transport, and incidence signs.  Hence the two
  ordered composites truncate to the same labelled composites in the lower
  window, and their pair cancellation survives.

  If a bracket row contains \(\alpha_{\mathrm{sc}}\) or \(\alpha_{11}\), one
  bracket input is zero as a local functional in the retained
  full-equivariant row family.  Bilinearity of the BV/convolution bracket
  gives \(R^{\mathrm{row}}_{b,M}=0\), and applying the scalar projection
  gives \(S_{b,M}=0\).  The zero truncation row is then compatible with the
  defining equations of
  Definition~\ref{defn:finite-window-graph-array}.

  For genuine seeds and for nonzero added bracket rows, the displayed
  formulas are not consequences of the cutoff alone; they are the required
  decompositions of the projected normalized row values in the chosen
  lower-window row basis.  Applying
  \(\pi_{N,P}\pi_{M,N}=\pi_{M,P}\) gives the displayed cocycle equations.
  Conversely, those equations are exactly the statement that the supplied
  matrices define a compatible finite-window row system.
\end{proof}

\begin{prop}[Projection formulas for the \(u\), \(v\), and \(q\) rows]\label{prop:projection-defined-uvq-truncation-data}
  Work in a larger finite-window graph package in the local
  \(\R^2_{\mathrm{top}}\times\C^2_{\mathrm{hol}}\) full-equivariant marked
  branch.  Fix \(M\geq N\).  Suppose the package supplies ordered
  lower-window row bases, or equivalently row presentations with chosen
  coordinate maps, for the following three finite-dimensional spans:
  \[
  \begin{aligned}
    \mathcal U_M
      &=
      \operatorname{span}\{\mathbf K^M_\Gamma\}_{\Gamma\in
      \mathcal G^{\mathrm{gen}}_{2,M}},
      \qquad
      \mathbf K^M_\Gamma=\varepsilon_\Gamma K^M_\Gamma,\\
    \mathcal V_M
      &=
      \operatorname{span}\{\mathbf E^M_{\Gamma,\nu}\}_{\Gamma,\nu},
      \qquad
      \mathbf E^M_{\Gamma,\nu}
        =
      -\varepsilon^{\mathrm{CE}}_{\Gamma,\nu}
      K^M_\Gamma(\zeta_{M,\nu}),\\
    \mathcal B_M
      &=
      \operatorname{span}\{\mathbf B^M_{\beta_1,\beta_2}\}_{\beta_i\notin
      \{\alpha_{\mathrm{sc}},\alpha_{11}\}},
      \qquad
      \mathbf B^M_{\beta_1,\beta_2}
        =
      \frac12\varepsilon^M_{\beta_1,\beta_2}
      \{K^M_{\beta_1},K^M_{\beta_2}\}_{\hbar,M}.
  \end{aligned}
  \]
  Assume these spans are projection-closed:
  \[
    \pi_{M,N}\mathcal U_M\subset\mathcal U_N,\qquad
    \pi_{M,N}\mathcal V_M\subset\mathcal V_N,\qquad
    \pi_{M,N}\mathcal B_M\subset\mathcal B_N,
  \]
  and that \(\pi_{M,N}\) is the chain and BV-bracket projection on the
  chosen finite-window subcomplex:
  \[
    \pi_{M,N}d_M=d_N\pi_{M,N},
    \qquad
    \pi_{M,N}\{x,y\}_{\hbar,M}
      =
    \{\pi_{M,N}x,\pi_{M,N}y\}_{\hbar,N}
  \]
  for the rows under consideration.

  Then the missing entries isolated above are defined by coordinates of
  finite-window projection:
  \[
    \pi_{M,N}\mathbf K^M_\Gamma
      =
    \sum_{\Gamma'}
      u^{M,N}_{\Gamma,\Gamma'}\mathbf K^N_{\Gamma'},
  \]
  \[
    \pi_{M,N}\mathbf E^M_{\Gamma,\nu}
      =
    \sum_{\Gamma',\nu'}
      v^{M,N}_{\Gamma,\nu;\Gamma',\nu'}
      \mathbf E^N_{\Gamma',\nu'},
  \]
  and
  \[
    \pi_{M,N}\mathbf B^M_{\beta_1,\beta_2}
      =
    \sum_{\beta'_1,\beta'_2}
      q^{M,N}_{\beta_1,\beta_2;\beta'_1,\beta'_2}
      \mathbf B^N_{\beta'_1,\beta'_2}.
  \]
  These equations are the definitions of \(u^{M,N}\), \(v^{M,N}\), and
  \(q^{M,N}\) relative to the chosen lower-window coordinates.  If the
  displayed lower-window families are honest bases, the entries are unique.
  If the package supplies only a spanning family with relations, the
  coordinates are an additional presentation datum; without it there is no
  canonical matrix entry.

  The \(u\)-entries also transport the differential rows:
  \[
    \pi_{M,N}(\varepsilon_\Gamma d_MK^M_\Gamma)
      =
    \sum_{\Gamma'}
      u^{M,N}_{\Gamma,\Gamma'}\,
      \varepsilon_{\Gamma'}d_NK^N_{\Gamma'}.
  \]
  This is just the first coordinate equation after applying
  \(d_N\pi_{M,N}=\pi_{M,N}d_M\).

  For the bracket rows there is a useful factorized form.  Suppose the
  added nonzero order-\(1\) rows have transport
  \[
    \pi_{M,N}K^M_\beta
      =
    \sum_{\eta}r^{M,N}_{\beta,\eta}K^N_\eta,
    \qquad
    \eta\notin\{\alpha_{\mathrm{sc}},\alpha_{11}\},
  \]
  and write the lower-window ordered raw bracket reduction as
  \[
    \frac12\{K^N_{\eta_1},K^N_{\eta_2}\}_{\hbar,N}
      =
    \sum_{\beta'_1,\beta'_2}
      \Theta^N_{\eta_1,\eta_2;\beta'_1,\beta'_2}
      \mathbf B^N_{\beta'_1,\beta'_2}.
  \]
  Then
  \[
    q^{M,N}_{\beta_1,\beta_2;\beta'_1,\beta'_2}
      =
    \varepsilon^M_{\beta_1,\beta_2}
    \sum_{\eta_1,\eta_2}
      r^{M,N}_{\beta_1,\eta_1}
      r^{M,N}_{\beta_2,\eta_2}
      \Theta^N_{\eta_1,\eta_2;\beta'_1,\beta'_2}.
  \]
  If the lower package retains all ordered normalized bracket rows and no
  further row relation, then
  \[
    \Theta^N_{\eta_1,\eta_2;\beta'_1,\beta'_2}
      =
    \frac{
      \delta_{\eta_1,\beta'_1}\delta_{\eta_2,\beta'_2}
    }{\varepsilon^N_{\beta'_1,\beta'_2}},
  \]
  and hence
  \[
    q^{M,N}_{\beta_1,\beta_2;\beta'_1,\beta'_2}
      =
    \frac{\varepsilon^M_{\beta_1,\beta_2}}
         {\varepsilon^N_{\beta'_1,\beta'_2}}\,
    r^{M,N}_{\beta_1,\beta'_1}
    r^{M,N}_{\beta_2,\beta'_2}.
  \]

  For \(M\geq N\geq P\), these coordinate matrices satisfy
  \[
    u^{M,P}_{\Gamma,\Gamma''}
      =
    \sum_{\Gamma'}
      u^{M,N}_{\Gamma,\Gamma'}u^{N,P}_{\Gamma',\Gamma''},
  \]
  \[
    v^{M,P}_{\Gamma,\nu;\Gamma'',\nu''}
      =
    \sum_{\Gamma',\nu'}
      v^{M,N}_{\Gamma,\nu;\Gamma',\nu'}
      v^{N,P}_{\Gamma',\nu';\Gamma'',\nu''},
  \]
  and
  \[
    q^{M,P}_{\beta_1,\beta_2;\theta_1,\theta_2}
      =
    \sum_{\beta'_1,\beta'_2}
      q^{M,N}_{\beta_1,\beta_2;\beta'_1,\beta'_2}
      q^{N,P}_{\beta'_1,\beta'_2;\theta_1,\theta_2}.
  \]

  Consequently the larger-package order-\(2\) non-scalar residual is a
  theorem-grade element of the compatible finite-window tower only after
  the following data are supplied: the genuine seed weights, source-face
  values, nonzero order-\(1\) rows, all orientation/Koszul/symmetry signs,
  CE-source transport, lower-window row bases or presentations, scalar-gate
  vanishing or scalar projections, and the projection-closure identities
  above.  With this package the obstruction is exactly the class of the
  resulting row sum in
  \(H^1(\mathcal K^\bullet_{\mathrm{ns},M}(I),d_{\mathrm{ns},M})\),
  together with the Roos compatibility class in
  \(\varprojlim^1 H^0(\mathcal K^\bullet_{\mathrm{ns},M}(I))\).  Without
  this package the non-scalar QME obstruction is not yet a defined class.
\end{prop}

\begin{proof}
  The first three displayed equations are simply the coordinate expansion of
  the projected normalized row values in the chosen lower-window row bases.
  Projection-closure is exactly the hypothesis that these coordinates
  exist.  If the lower rows are linearly independent, the coordinate
  expansion is unique.  If there are row relations, uniqueness is lost unless
  the package supplies a presentation or a coordinate splitting.

  The \(d\Gamma\)-row formula follows by applying the chain-map identity to
  \(\pi_{M,N}\mathbf K^M_\Gamma\).  The \(q\)-formula follows by applying the
  BV-bracket identity:
  \[
  \begin{aligned}
    \pi_{M,N}\mathbf B^M_{\beta_1,\beta_2}
      &=
      \frac12\varepsilon^M_{\beta_1,\beta_2}
      \{\pi_{M,N}K^M_{\beta_1},
        \pi_{M,N}K^M_{\beta_2}\}_{\hbar,N}  \\
      &=
      \frac12\varepsilon^M_{\beta_1,\beta_2}
      \sum_{\eta_1,\eta_2}
      r^{M,N}_{\beta_1,\eta_1}
      r^{M,N}_{\beta_2,\eta_2}
      \{K^N_{\eta_1},K^N_{\eta_2}\}_{\hbar,N},
  \end{aligned}
  \]
  and then reducing each lower raw bracket by the matrix \(\Theta^N\).
  The ordered normalized-row specialization is the case where the lower
  raw bracket is already one basis row after division by its incidence
  coefficient.

  The cocycle identities are forced by
  \(\pi_{N,P}\pi_{M,N}=\pi_{M,P}\).  Expanding first directly from \(M\) to
  \(P\), and then through \(N\), gives two coordinate expansions of the same
  projected row in the same lower-window coordinates.  Equality of
  coordinates gives the three displayed matrix product formulas.  These
  compatibility identities are precisely the truncation part of
  Definition~\ref{defn:finite-window-graph-array}.  Once they are supplied,
  the obstruction criterion is the non-scalar kernel criterion of
  Definition~\ref{def:app-nonscalar-kernel-row-complex} and
  Proposition~\ref{prop:finite-window-qme-limit-condition}.
\end{proof}

\begin{prop}[Finite-row primitive search interface]\label{prop:finite-row-primitive-search-interface}
  Let \(\mathsf A^{\mathrm{fw}}_{n,M}\) be a valid finite-window graph
  array at order \(n\), and suppose its scalar gate has vanished:
  \[
    \sum_{\alpha\in\mathsf A^{\mathrm{fw}}_{n,M}}S_{\alpha,M}=0.
  \]
  Choose an ordered row presentation
  \[
    \rho^M_1,\ldots,\rho^M_a
  \]
  for the degree-one non-scalar rows retained at this order: graph
  differentials, CE-source faces, nonzero bracket rows, and the lower
  counterterm rows.  Write the residual as
  \[
    R_{n,M}=\sum_i r^M_i\rho^M_i.
  \]
  Choose also an ordered finite list of degree-zero counterterm candidates
  \[
    \eta^M_1,\ldots,\eta^M_b\in
    \mathcal K^0_{\mathrm{ns},M}(I)
  \]
  supplied by the same finite package, and write
  \[
    d_{\mathrm{ns},M}\eta^M_j=\sum_i A^M_{ij}\rho^M_i.
  \]
  Then a counterterm in the supplied candidate span,
  \(C_{n,M}=\sum_j c_j\eta^M_j\), exists if and only if the finite linear
  system
  \[
    A^M c=-r^M
  \]
  is solvable.  If it is not solvable, any covector
  \(\ell\in{\left(\C^a\right)}^\vee\) satisfying
  \[
    \ell A^M=0,\qquad \ell(r^M)\neq0
  \]
  is a certificate that the class of \(R_{n,M}\) is nonzero in the
  cohomology of the displayed finite row complex.

  For a tower \(M\geq N\), let \(T^{M,N}\) be the row-truncation matrix on
  the \(\rho\)-presentation and \(P^{M,N}\) the chosen transport matrix on
  the primitive candidates.  The row block \(T^{M,N}\) is the direct sum of
  the already computed identity or zero blocks, the \(u\)-block for genuine
  seed rows, the \(v\)-block for CE-source faces, and the \(q\)-block for
  nonzero bracket rows.  A compatible primitive search is defined only when
  \[
    T^{M,N}r^M=r^N,\qquad
    T^{M,N}A^M=A^N P^{M,N}.
  \]
  After windowwise solutions \(c_M\) have been chosen, the remaining
  compatibility obstruction is the Roos class represented by
  \[
    \bigl[P^{M,N}c_M-c_N\bigr]\in
    H^0(\mathcal K^\bullet_{\mathrm{ns},N}(I)).
  \]
  Thus the primitive search has three gates, in this order: scalar gate,
  fixed-window linear solvability, and Roos \(\varprojlim^1\)
  compatibility.  The last gate is not reached when a fixed-window
  \(H^1\)-obstruction is already certified by such a covector \(\ell\).
\end{prop}

\begin{proof}
  Once the scalar gate vanishes, the residual lies in the non-scalar kernel
  complex.  Restricting the possible counterterm to the supplied degree-zero
  span turns the equation
  \(d_{\mathrm{ns},M}C_{n,M}=-R_{n,M}\) into the displayed matrix equation.
  Failure of solvability is equivalent to \(-r^M\notin\operatorname{im}A^M\).
  In finite dimension this is certified by a linear functional annihilating
  \(\operatorname{im}A^M\) and not annihilating \(r^M\), exactly as written.

  The truncation identities say that the row residual and the boundaries of
  candidate primitives commute with finite-window projection.  The computed
  blocks are those of
  Propositions~\ref{prop:computed-finite-window-truncation-matrices}
  and~\ref{prop:projection-defined-uvq-truncation-data}.  If \(c_M\) and \(c_N\)
  solve the two fixed-window equations, then
  \(P^{M,N}c_M-c_N\) is closed in degree zero by
  \(T^{M,N}A^M=A^N P^{M,N}\) and \(T^{M,N}r^M=r^N\).  Its cohomology class is
  the Roos \(1\)-cocycle of
  Proposition~\ref{prop:finite-window-qme-limit-condition}.
\end{proof}

\begin{prop}[Concrete order-three larger-package row]\label{prop:concrete-order-three-larger-package-row}
  Work in the full-equivariant marked branch with \(V=\C^{N|N}\), and let
  \(M\geq3\).  Enlarge the minimal finite-window package by one Hom-valued
  CE-source face row
  \[
    \theta_3=\mathrm{CE}(\Theta_3,\nu_3),
    \qquad |\Theta_3|_\hbar=3,
  \]
  where \(\Theta_3\) is a genuine marked Costello seed, not a
  scalar-contact closure composite and not a bracket row.  The row data are
  as follows.

  The graph type and order are
  \[
    \tau_{\theta_3}=\mathrm{CE},\qquad |\theta_3|_\hbar=3.
  \]
  Its finite labels are
  \[
  \begin{aligned}
    e_{\mathrm{pp}}&:\quad
      P_{\epsilon,L}^{w,M},\quad
      h_{2,1}=z_1^2z_2,\ \rho_{2,1},\\
    e_{\mathrm{ct}}&:\quad
      \Delta_{\epsilon,L}^{w,M},\quad
      h_{1,1}=z_1z_2,\ \rho_{1,1},\\
    e_{\mathrm{W}}&:\quad
      E_{\mathrm{Weyl}},\quad
      z_1^2,\ z_2
  \end{aligned}
  \]
  with \(B_\theta\in\mathcal D^{\mathrm{reg}}_c(I)\) on the cotangent
  external leg.  The vertex labels are
  \[
    u_{a(t)dt\otimes z_1^2},\qquad
    u_{b(t)dt\otimes z_2},\qquad
    c_{B_\theta,\rho_{2,1}},\qquad
    I^w_{0,M},\qquad
    W^{R,M}_{\Theta_3,L,w}(B^\Theta_\bullet).
  \]
  The marker tensor is a full
  \(\lie{gl}(N|N)\)-equivariant signed Brauer marker
  \[
    \mathfrak m_{\theta_3}\in
    \mathcal M^{\mathrm{full}}(\C^{N|N}),
    \qquad
    \operatorname{Str}_{\mathrm{cyc}}(\mathfrak m_{\theta_3})=0.
  \]
  Choose the source-face coefficient
  \[
    d_{\mathrm{CE}}\zeta_M
      =
    +\zeta_{M,\nu_3}+\text{other faces}
  \]
  on the \(\Theta_3\)-component.  Since the kernel has degree zero, the
  incidence sign of this CE row in
  \(-\kappa^M_{\mathcal G,w,I}(d_{\mathrm{CE}}\zeta_M)\) is \(-1\).

  Write
  \[
    K^M_{\Theta_3}
      =
    W^{R,M}_{\Theta_3,L,w}(B^\Theta_\bullet).
  \]
  The row value is the nonzero symbolic local functional
  \[
    R^{\mathrm{row}}_{\theta_3,M}
      =
    -K^M_{\Theta_3}(\zeta_{M,\nu_3})
      =:\mathsf E_{\theta_3,M}\neq0.
  \]
  Its scalar gate is
  \[
    S_{\theta_3,M}
      =
    \widehat\sigma_{\mathrm{sc},M}(\mathsf E_{\theta_3,M})=0.
  \]
  Indeed, the row lies in the full-equivariant marked habitat, so
  Theorem~\ref{thm:app-full-equivariant-marked-shadow-vanishing} kills its
  scalar projection; also the displayed Weyl half-edge has
  \(\omega(z_1^2,z_2)=0\).

  The class-forming closure check is part of the row datum:
  \[
    d_{\mathrm{ns},M}\mathsf E_{\theta_3,M}=0.
  \]
  If this equation is not supplied, the row is not a valid obstruction row
  and no \(H^1\)-class is formed.  With it supplied, the row defines
  \[
    \mathfrak o^{\mathrm{ns}}_{\theta_3,M}
      =
    [\mathsf E_{\theta_3,M}]
      \in
    H^1\!\left(
      \mathcal K^\bullet_{\mathrm{ns},M}(I),d_{\mathrm{ns},M}
    \right).
  \]

  The truncation matrix is explicit.  For \(M\geq N\),
  \[
    \pi_{M,N}\mathsf E_{\theta_3,M}
      =
    \begin{cases}
      \mathsf E_{\theta_3,N}, & N\geq3,\\
      0, & N<3,
    \end{cases}
  \]
  hence
  \[
    t^{M,N}_{\theta_3\theta_3}=1\quad(N\geq3),\qquad
    t^{M,N}_{\theta_3\alpha'}=0\quad(\alpha'\neq\theta_3),
  \]
  and every entry from \(\theta_3\) is zero for \(N<3\).  These entries
  satisfy \(t^{M,P}=t^{M,N}t^{N,P}\).

  Therefore this larger-package row does not vanish as a local-functional
  cochain, but its scalar gate vanishes.  It is an explicit symbolic
  obstruction datum.  A fixed-window order-three counterterm exists
  precisely when
  \[
    \mathfrak o^{\mathrm{ns}}_{\theta_3,M}=0,
    \qquad
    d_M C_{\theta_3,M}=-\mathsf E_{\theta_3,M}
  \]
  for some \(C_{\theta_3,M}\in\mathcal K^0_{\mathrm{ns},M}(I)\).  No such
  primitive is supplied by this row, so the manuscript may name the class
  but may not claim the larger order-three QME branch is solved.
\end{prop}

\begin{prop}[\(\theta_3\) finite-row obstruction certificate]\label{prop:theta-three-finite-row-obstruction}
  Keep the one-row order-three package of
  Proposition~\ref{prop:concrete-order-three-larger-package-row}.  In the
  finite row complex actually supplied there, set
  \[
    \mathcal K^0_{\theta_3,M}=0,\qquad
    \mathcal K^1_{\theta_3,M}=\C\mathsf E_{\theta_3,M},\qquad
    d_{\theta_3,M}=0,
  \]
  and take residual \(R_{\theta_3,M}=\mathsf E_{\theta_3,M}\).  The scalar
  gate has already vanished:
  \[
    \widehat\sigma_{\mathrm{sc},M}(\mathsf E_{\theta_3,M})=0.
  \]
  The primitive-search matrix is the \(1\times0\) zero matrix and the target
  vector is \(-\mathsf E_{\theta_3,M}\).  Hence the equation
  \(A^M c=-r^M\) has no solution in the displayed finite row complex.
  The covector
  \[
    \ell_{\theta_3}(\mathsf E_{\theta_3,M})=1
  \]
  annihilates \(\operatorname{im}d_{\theta_3,M}=0\) and satisfies
  \(\ell_{\theta_3}(R_{\theta_3,M})=1\).  Therefore
  \[
    [\mathsf E_{\theta_3,M}]\neq0
    \quad\text{in}\quad
    H^1(\mathcal K^\bullet_{\theta_3,M}).
  \]
  This is an exact obstruction statement for the finite row subcomplex
  presently supplied.  It does not exclude a primitive in a larger
  local-functional complex.  To heal the row one must add either a
  degree-zero candidate \(C_{\theta_3,M}\) with
  \(d_{\mathrm{ns},M}C_{\theta_3,M}=-\mathsf E_{\theta_3,M}\), or additional
  Hom-valued companion rows whose sum changes the residual vector.

  If such a candidate \(C_{\theta_3,M}\) is supplied and truncates by
  \[
    \pi_{M,N}C_{\theta_3,M}=C_{\theta_3,N}\quad(N\geq3),
    \qquad
    \pi_{M,N}C_{\theta_3,M}=0\quad(N<3),
  \]
  then the fixed-window search is solved and the Roos
  \(\varprojlim^1\)-class of this one-row tower is zero.  If the
  truncation identity fails, the residual fixed-window obstruction is gone
  but the Roos class is represented by
  \[
    [\pi_{M,N}C_{\theta_3,M}-C_{\theta_3,N}]
    \in H^0(\mathcal K^\bullet_{\mathrm{ns},N}(I)).
  \]
  For the minimal order-two full-equivariant rows and for bracket rows
  containing \(\alpha_{\mathrm{sc}}\) or \(\alpha_{11}\), the same primitive
  search has residual vector \(0\), so the zero counterterm is the unique
  supplied solution and its Roos class is zero.
\end{prop}

\begin{proof}
  This is Proposition~\ref{prop:finite-row-primitive-search-interface}
  applied to the one-row presentation.  Since no degree-zero row is supplied,
  the image of the finite differential is zero.  The residual vector is the
  basis vector \(\mathsf E_{\theta_3,M}\), so the displayed covector is a
  cokernel certificate.  The Roos statement is the same proposition applied
  after a hypothetical primitive has been added.  The minimal order-two and
  zero-bracket rows have zero residual before cohomology, hence the zero
  primitive is compatible under truncation.
\end{proof}

\begin{prop}[\(\theta_3\) payload acceptance gate]\label{prop:theta-three-payload-acceptance-gate}
  Keep the row \(\theta_3=\mathrm{CE}(\Theta_3,\nu_3)\) and the notation of
  Propositions~\ref{prop:concrete-order-three-larger-package-row}
  and~\ref{prop:theta-three-finite-row-obstruction}.  A proposed
  degree-zero healing payload for this row is accepted by the finite-row
  theorem only after the row array supplies one of the following two
  primitive entries.
  \begin{enumerate}[label=(P\arabic*)]
    \item A CE-ancestor entry
      \[
        \eta_{\theta_3,M},\qquad
        d_{\mathrm{CE}}\eta_{\theta_3,M}=\zeta_{M,\nu_3},
      \]
      together with the Hom-valued Bianchi and companion-face terms needed
      to prove
      \[
        C^{\mathrm{CE}}_{\theta_3,M}
          =
        K^M_{\Theta_3}(\eta_{\theta_3,M})
        \in\mathcal K^0_{\mathrm{ns},M}(I),
        \qquad
        d_{\mathrm{ns},M}C^{\mathrm{CE}}_{\theta_3,M}
          =
        -\mathsf E_{\theta_3,M}.
      \]

    \item A local counterterm entry
      \(C^{\mathrm{ct}}_{\theta_3,M}\in
      \mathcal K^0_{\mathrm{ns},M}(I)\), with its locality proof, scalar
      projection, and differential expansion, satisfying
      \[
        d_{\mathrm{ns},M}C^{\mathrm{ct}}_{\theta_3,M}
          =
        -\mathsf E_{\theta_3,M}.
      \]
  \end{enumerate}
  Equivalently, in the ordered one-row presentation
  \[
    \rho^M_1=\mathsf E_{\theta_3,M},\qquad r^M=(1),
  \]
  the supplied candidate column \(C_{\theta_3,M}\) must have boundary-matrix
  coefficient
  \[
    A^M_{\theta_3,C}=-1,
    \qquad
    d_{\mathrm{ns},M}C_{\theta_3,M}
      =
    -\mathsf E_{\theta_3,M}.
  \]
  A row-data payload with no displayed differential entry, or with any
  coefficient other than \(-1\) for the named candidate, is not this
  accepted payload; it must first be replaced by an explicitly supplied
  candidate whose differential column is the one above.

  The scalar and tower gates are part of the same datum:
  \[
    \widehat\sigma_{\mathrm{sc},M}(C_{\theta_3,M})=0,
    \qquad
    \pi_{M,N}C_{\theta_3,M}=
    \begin{cases}
      C_{\theta_3,N}, & N\geq3,\\
      0, & N<3.
    \end{cases}
  \]
  With the row transport
  \(T^{M,N}_{\theta_3\theta_3}=1\) for \(N\geq3\) and \(0\) for
  \(N<3\), these equations give
  \[
    T^{M,N}A^M=A^N P^{M,N},
    \qquad
    [P^{M,N}c_M-c_N]=0
    \quad\text{for the solution }c_M=1.
  \]
  Thus such a payload kills the fixed-window row and has zero Roos
  compatibility class.  The present manuscript data do not contain this
  payload: in the supplied one-row complex \(K^0_{\theta_3,M}=0\), so
  Proposition~\ref{prop:theta-three-finite-row-obstruction} remains the
  active claim.

  A complete companion-face table is a different gate.  It does not supply a
  primitive.  Instead it replaces the one-face residual by a full
  CE-source-face row sum
  \[
    \sum_{\nu\in F_{\Theta_3,M}}
      R^{\mathrm{row}}_{\mathrm{CE}(\Theta_\nu,\nu),M},
    \qquad
    \nu_3\in F_{\Theta_3,M},
  \]
  whose \(\nu_3\)-summand is \(\mathsf E_{\theta_3,M}\).  Such a table is
  accepted only if it supplies every face, every CE coefficient, every
  scalar gate, the \(v^{M,N}\)-truncation matrices with their cocycle
  identity, and the codimension-two closure table, and if
  \[
    \sum_{\nu\in F_{\Theta_3,M}}
      R^{\mathrm{row}}_{\mathrm{CE}(\Theta_\nu,\nu),M}=0,
    \qquad
    \sum_{\nu\in F_{\Theta_3,M}}
      S_{\mathrm{CE}(\Theta_\nu,\nu),M}=0.
  \]
  The current one-face table is only an interface for this possibility: its
  residual is \(\mathsf E_{\theta_3,M}\), not zero.  Hence it does not solve
  the \(\theta_3\) row and does not form a primitive Roos class.  The first
  source-algebraic companion table containing \(\nu_3\) is recorded in
  Theorem~\ref{thm:theta-three-companion-face-obstruction}; it produces an
  obstruction datum, not a cancellation theorem.
\end{prop}

\begin{proof}
  For a primitive payload, apply
  Proposition~\ref{prop:finite-row-primitive-search-interface} with
  degree-one basis \(\rho^M_1=\mathsf E_{\theta_3,M}\) and residual vector
  \(r^M=(1)\).  The equation \(A^M c=-r^M\) is solved by the supplied
  candidate with \(c=1\) exactly when the candidate column is
  \(A^M_{\theta_3,C}=-1\).  The CE-ancestor and local-counterterm
  alternatives are precisely two ways of producing such a degree-zero
  column; without the displayed target differential, scalar projection, and
  support proof, the finite row complex has no boundary matrix entry to use.

  The truncation assertion is the identity/zero transport already attached
  to the row \(\mathsf E_{\theta_3,M}\).  With the same identity/zero rule
  for \(C_{\theta_3,M}\), the matrix equation \(T A=A P\) is tautological:
  above degree \(3\) it reads \((-1)=(-1)\), and below degree \(3\) both
  sides are \(0\).  The chosen coefficients \(c_M=1\) are therefore
  compatible, so the Roos representative vanishes.

  For a companion-face table no degree-zero candidate is chosen.  The
  residual vector itself is replaced by the signed row sum of all CE-source
  faces.  The finite-row theorem can test that new row sum only after the
  table supplies its scalar gates, truncation coordinates, and
  codimension-two closure data.  If the signed sum is not zero, the row
  remains a residual; if it is zero, cancellation occurs before the
  primitive search and no primitive Roos class is formed.
\end{proof}

\begin{thm}[\(\theta_3\) companion-face construction obstruction]\label{thm:theta-three-companion-face-obstruction}
  Work in the order-three full-equivariant marked branch of
  Proposition~\ref{prop:concrete-order-three-larger-package-row}, and put
  \(H_{p,q}=z_1^pz_2^q\).  Consider the two-\(u\) source-algebra ancestor
  \[
    \zeta^0_M=u_{a,H_{1,0}}u_{b,H_{0,1}},
    \qquad
    [H_{1,0},H_{2,1}]=H_{2,0}.
  \]
  Its monomial CE differential is the first source-algebraic table of the
  same external shape whose source faces contain the displayed
  \(\nu_3\)-row
  \(u_{a,H_{2,0}}u_{b,H_{0,1}}c_{B_\theta,\rho_{2,1}}\).  After truncation
  to Hamiltonian degree \(\leq3\), removal of constants, and cancellation of
  the two \(\rho_{1,1}\)-terms, the nonzero candidate face set is the
  following eight-face set.  With
  \[
    E^M_\nu=-K^M_{\Theta_3}(\zeta_{M,\nu}),
    \qquad
    E^M_{\nu_3}=\mathsf E_{\theta_3,M},
  \]
  the table is
  \[
  \begin{array}{c|c|l|l}
    \nu & \epsilon^{\mathrm{CE}}_{\Theta_3,\nu}
      & \zeta_{M,\nu} & \epsilon^{\mathrm{CE}}_{\Theta_3,\nu}E^M_\nu \\
    \hline
    \nu_{02} & +2
      & u_{a,H_{0,1}}u_{b,H_{0,1}}c_{B_\theta,\rho_{0,2}}
      & 2E^M_{\nu_{02}} \\
    \nu_{20} & -2
      & u_{a,H_{1,0}}u_{b,H_{1,0}}c_{B_\theta,\rho_{2,0}}
      & -2E^M_{\nu_{20}} \\
    \nu_{03} & +3
      & u_{a,H_{0,2}}u_{b,H_{0,1}}c_{B_\theta,\rho_{0,3}}
      & 3E^M_{\nu_{03}} \\
    \nu_{12}^{(1)} & +2
      & u_{a,H_{1,1}}u_{b,H_{0,1}}c_{B_\theta,\rho_{1,2}}
      & 2E^M_{\nu_{12}^{(1)}} \\
    \nu_{12}^{(2)} & -1
      & u_{a,H_{1,0}}u_{b,H_{0,2}}c_{B_\theta,\rho_{1,2}}
      & -E^M_{\nu_{12}^{(2)}} \\
    \nu_3=\nu_{21}^{(1)} & +1
      & u_{a,H_{2,0}}u_{b,H_{0,1}}c_{B_\theta,\rho_{2,1}}
      & \mathsf E_{\theta_3,M} \\
    \nu_{21}^{(2)} & -2
      & u_{a,H_{1,0}}u_{b,H_{1,1}}c_{B_\theta,\rho_{2,1}}
      & -2E^M_{\nu_{21}^{(2)}} \\
    \nu_{30} & -3
      & u_{a,H_{1,0}}u_{b,H_{2,0}}c_{B_\theta,\rho_{3,0}}
      & -3E^M_{\nu_{30}}
  \end{array}
  \]
  Hence the source-algebraic companion residual is
  \[
  \begin{aligned}
    R^{\mathrm{cand}}_{\Theta_3,M}
      ={}&2E^M_{\nu_{02}}-2E^M_{\nu_{20}}
          +3E^M_{\nu_{03}}+2E^M_{\nu_{12}^{(1)}}
          -E^M_{\nu_{12}^{(2)}}\\
       &+\mathsf E_{\theta_3,M}
          -2E^M_{\nu_{21}^{(2)}}-3E^M_{\nu_{30}} .
  \end{aligned}
  \]
  This expression is the exact residual of the attempted companion route.
  It is not known to be zero.

  The induced source-face transport is the diagonal cutoff
  \[
    v^{M,N}_{\nu;\nu'}
      =
    \begin{cases}
      1,&\nu=\nu'\ \text{and}\ d(\nu)\leq N,\\
      0,&\text{otherwise},
    \end{cases}
  \]
  where \(d(\nu_{02})=d(\nu_{20})=2\) and the other six faces have
  degree \(3\).  It satisfies the cocycle identity
  \(v^{M,P}=v^{M,N}v^{N,P}\).  For \(N=2\), however, the transported
  residual is
  \[
    \pi_{M,2}R^{\mathrm{cand}}_{\Theta_3,M}
      =
    2E^2_{\nu_{02}}-2E^2_{\nu_{20}}.
  \]
  Thus even a fixed-window cancellation at \(M\geq3\) would not give a
  theorem-grade compatible tower unless the lower-window identity
  \(E^2_{\nu_{02}}=E^2_{\nu_{20}}\) is proved in the chosen row basis, or
  the lower-window row presentation is changed and its \(v\)-matrix is
  recomputed.

  Consequently the companion route is a construction/obstruction theorem
  with the following exact missing marked Costello datum:
  \[
    \left(
      d_{\mathrm{CE}}\zeta_M,\,
      \epsilon^{\mathrm{CE}}_{\Theta_3,\nu},\,
      K^M_{\Theta_3}(\zeta_{M,\nu}),\,
      \mu_\nu,\,
      v^{M,N}_{\nu;\nu'},\,
      \mathcal C^{(2)}_{\nu,\nu'}
    \right)_{\nu,\nu'} .
  \]
  Here \(\mu_\nu\) is the marked output and marker-transport entry, and
  \(\mathcal C^{(2)}_{\nu,\nu'}\) is the codimension-two source-face
  incidence/weight table: the two orders of each pair of boundary
  operations, the output graph, the orientation/Koszul coefficient, the
  transported marker, and the resulting renormalized row coordinate in the
  declared degree-one basis.  Until this table is supplied and proves both
  \(R^{\mathrm{cand}}_{\Theta_3,M}=0\) and its transported lower-window
  identities, the current theorem surface remains the one-row obstruction
  of Proposition~\ref{prop:theta-three-finite-row-obstruction}.  No
  companion cancellation and no primitive counterterm is claimed.
\end{thm}

\begin{proof}
  The Poisson rule
  \([H_{a,b},H_{r,s}]=(as-br)H_{a+r-1,b+s-1}\) gives the displayed
  source-algebra faces from the differential of the two \(u\)-factors
  against the current label.  In degree \(\leq3\), constants are removed by
  the Hamiltonian quotient and the two \(\rho_{1,1}\)-faces have opposite
  coefficients, so the eight listed terms remain.  Applying the row
  convention
  \(R^{\mathrm{row}}_{\mathrm{CE}(\Theta_3,\nu),M}
  =\epsilon^{\mathrm{CE}}_{\Theta_3,\nu}E^M_\nu\) gives the signed residual
  formula, with the \(\nu_3\)-summand equal to
  \(\mathsf E_{\theta_3,M}\).

  The diagonal transport is just finite-window projection on the Hamiltonian
  degree of the \(c\)-label.  Such diagonal cutoffs compose, proving the
  cocycle identity.  Projecting to \(N=2\) deletes every degree-three face
  and leaves exactly the two degree-two summands.  Therefore the lower-window
  residual is \(2E^2_{\nu_{02}}-2E^2_{\nu_{20}}\).

  The passage from this source-algebra table to a Costello cancellation
  theorem is precisely the missing marked incidence/weight computation.
  Source coefficients alone do not determine the renormalized graph weights,
  marker transports, codimension-two closure signs, or lower-window row
  coordinates.  Those entries are the displayed missing datum.  Without
  them, Definition~\ref{defn:finite-window-graph-array} supplies no zero
  residual vector, and Proposition~\ref{prop:finite-row-primitive-search-interface}
  still sees the current one-row obstruction.
\end{proof}

\begin{thm}[Actual finite-window non-scalar decision datum]\label{thm:actual-finite-window-nonscalar-decision-datum}
  Fix a finite-window graph/QME package
  \((I,w,M,\mathcal G_M)\), and fix the labelled CE cycle
  \(\zeta_M\) on which the curvature is evaluated.  Write
  \(K^M_{\Gamma}\) for the component of
  \(\kappa^M_{\mathcal G,w,I}\) attached to a graph
  \(\Gamma\); for the zero-internal-edge component this is
  \(\kappa^{\mathrm{coef},M}_{\hbar,w,I}\), and for a positive graph it
  is the renormalized weight \(W^{R,M}_{\Gamma,L,w}\).  If
  \(|\Gamma|_\hbar\) denotes the \(\hbar\)-order assigned by the finite
  graph package, then the order-\(n\) graph curvature is
  \[
  \begin{aligned}
    \operatorname{Ob}^{\mathrm{graph}}_{\mathcal G,M}[n](\zeta_M)
      ={}&[\hbar^n]\Bigl(
        d_M\kappa^M_{\mathcal G,w,I}(\zeta_M)
        -\kappa^M_{\mathcal G,w,I}(d_{\mathrm{CE}}\zeta_M)\\
      &\hspace{3.2cm}
        +\frac12
        \sum_{(\zeta_M)}
        \{\kappa^M_{\mathcal G,w,I}(\zeta'_M),
          \kappa^M_{\mathcal G,w,I}(\zeta''_M)\}_{\hbar,M}
      \Bigr)                                                     \\
      ={}&
      \sum_{|\Gamma|_\hbar=n}
        \varepsilon_\Gamma\,d_M K^M_\Gamma
      -\sum_{|\Gamma'|_\hbar=n}
        \varepsilon^{\mathrm{CE}}_{\Gamma'}K^M_{\Gamma'}          \\
      &\quad
      +\frac12
      \sum_{|\Gamma_1|_\hbar+|\Gamma_2|_\hbar=n}
        \varepsilon_{\Gamma_1,\Gamma_2}
        \{K^M_{\Gamma_1},K^M_{\Gamma_2}\}_{\hbar,M}.
  \end{aligned}
  \]
  The signs are the orientation signs of the supplied finite graph
  differential and the CE coproduct.  With lower counterterms
  \(C_{<n,M}\), the cochain tested at order \(n\) is
  \[
    R_{n,M}
      =
    \operatorname{Ob}^{\mathrm{graph}}_{\mathcal G,M}[n]
      +[\hbar^n]\left(
        d_M C_{<n,M}
        +\frac12\{C_{<n,M},C_{<n,M}\}_{\hbar,M}\right).
  \]
  Its scalar gate is
  \[
    \mathfrak s_{n,M}=\widehat\sigma_{\mathrm{sc},M}(R_{n,M}).
  \]

  The first scalar-contact graph whose coefficient is actually computed
  in the present local model is the one-cross-contraction boundary graph
  with two Hamiltonian external legs and one open scalar loop.  It has
  \(\hbar\)-order \(1\).  For a full-equivariant marked boundary
  generator \(X_{\Gamma,o,\mathfrak m,M}\), the extracted order-one
  scalar gate is
  \[
    \mathfrak s^{\mathrm{mk}}_{1,M}
      (X_{\Gamma,o,\mathfrak m,M})
      =
    \sum_{i:\,r(\beta_i)=1}
      (-1)^{i-1}
      \operatorname{Str}_{\mathrm{cyc}}
        \bigl(\mu_{\beta_i}(\mathfrak m)\bigr)\,
      \omega(f_{\partial_{\beta_i}\Gamma},
             g_{\partial_{\beta_i}\Gamma})
      \gamma_{\mathbf 1}.
  \]
  Before the coefficient extraction this is the
  \(\hbar^1\)-term
  \(\hbar\,\mathfrak s^{\mathrm{mk}}_{1,M}\).
  If \(V=\C^{N|N}\) and the markers are full
  \(\lie{gl}(V)\)-equivariant, then
  \[
    \operatorname{Str}_{\mathrm{cyc}}
      \bigl(\mu_{\beta_i}(\mathfrak m)\bigr)=0
  \]
  for every elementary boundary term by
  Lemma~\ref{lem:app-full-equivariant-cyclic-supertrace}; hence
  \[
    \mathfrak s^{\mathrm{mk}}_{1,M}=0.
  \]
  This is a scalar-gate vanishing, not a construction of a non-scalar
  primitive.

  After the scalar gate vanishes, the first non-scalar order is the least
  \(n_0\) for which the cochain \(R_{n_0,M}\) is defined for the actual
  graph list and is not already killed by the lower-order choices.  Its
  class is
  \[
    [\mathfrak o^{\mathrm{fw}}_{n_0,M}]
      =
    [R_{n_0,M}]
      \in
    H^1\bigl(\mathcal Q^\bullet_{\mathcal G,M}(I),d_M\bigr).
  \]
  It vanishes exactly when there exists
  \(C_{n_0,M}\in\mathcal Q^0_{\mathcal G,M}(I)\) with
  \(d_M C_{n_0,M}=-R_{n_0,M}\).  In a compatible finite-window tower,
  after choosing any such windowwise primitives, the remaining
  compatibility obstruction is
  \[
    \lambda^{\mathrm{fw}}_{n_0}
      =
    \left[
      \bigl(
        [\pi_{M+1,M}C_{n_0,M+1}-C_{n_0,M}]
      \bigr)_M
    \right]
      \in
    \varprojlim\nolimits^1_M
    H^0\bigl(\mathcal Q^\bullet_{\mathcal G,M}(I)\bigr).
  \]
  Thus compatible finite-window counterterms exist precisely when all
  fixed-window classes \([\mathfrak o^{\mathrm{fw}}_{n_0,M}]\) vanish and
  \(\lambda^{\mathrm{fw}}_{n_0}=0\).

  The data above are not recoverable from the abstract existence of
  \(\mathcal G_M\).  To decide \(n_0\), the class
  \([\mathfrak o^{\mathrm{fw}}_{n_0,M}]\), or the Roos class
  \(\lambda^{\mathrm{fw}}_{n_0}\), one must supply a valid finite-window
  coefficient array in the sense of
  Definition~\ref{defn:finite-window-graph-array}:
  \[
    \mathsf A^{\mathrm{fw}}_{n,M}
  \]
  whose rows specify graph type, edge labels, vertex labels, marker tensor,
  incidence coefficient, \(\hbar\)-order, coefficient degree,
  scalar-contact filtration, scalar-gate term, row curvature value, and
  truncation coefficients in \(M\).  Without this array the manuscript has
  only the formal obstruction symbol of
  Theorem~\ref{thm:finite-window-graph-qme-assembly}; it has no
  theorem-grade computation of \(n_0\), no proof that
  \([\mathfrak o^{\mathrm{fw}}_{n_0,M}]=0\), and no defined
  \(\lambda^{\mathrm{fw}}_{n_0}\).
\end{thm}

\begin{proof}
  The first displayed formula is the definition of the convolution
  curvature in
  Definition~\ref{def:app-bulk-defect-convolution-habitat}, evaluated on
  \(\zeta_M\) and then projected to the \(\hbar^n\)-coefficient.  Expanding
  \(-\kappa d_{\mathrm{CE}}\) and the CE coproduct bracket gives the
  three graph sums, with exactly the orientation signs supplied by the
  finite graph package.  Adding the lower counterterm terms gives
  \(R_{n,M}\), the residual of
  Theorem~\ref{thm:finite-window-graph-qme-assembly}.

  The order-one scalar-contact computation is
  Proposition~\ref{prop:app-first-nonidentity-loop-scalar-symbol} applied
  to each first filtration-lowering marked boundary term of
  Definition~\ref{def:app-oriented-full-equivariant-marked-boundary-complex}.
  The full-equivariant marker transport maps preserve the signed Brauer
  marker span.  Lemma~\ref{lem:app-full-equivariant-cyclic-supertrace}
  computes the cyclic supertrace of every such output marker as a power
  of \(\operatorname{sdim}\C^{N|N}=0\), and the empty marker word gives the
  same zero coefficient.  Therefore the scalar gate vanishes on the
  full-equivariant marked habitat before any non-scalar cohomology class
  is formed.

  Once the scalar gate is zero, \(R_{n,M}\) lies in
  \(\mathcal Q^\bullet_{\mathcal G,M}(I)=
  \ker\widehat\sigma_{\mathrm{sc},M}\) exactly as in the proof of
  Theorem~\ref{thm:finite-window-graph-qme-assembly}.  A degree-zero
  counterterm changes \(R_{n,M}\) by a \(d_M\)-boundary, so the
  fixed-window solvability criterion is precisely the vanishing of its
  \(H^1\)-class.  If primitives are chosen window by window, their
  truncation differences are closed degree-zero cochains; their classes
  form the standard Roos cocycle computing
  \(\varprojlim^1_M H^0(\mathcal Q^\bullet_{\mathcal G,M}(I))\).  Vanishing
  of this class is exactly the condition that the primitives can be
  adjusted to a compatible family.

  The final assertion is a dependency statement.  The formulas require a
  finite graph list with order assignment, CE cycle, signs, coproduct
  decomposition, graph weights, lower counterterms, scalar projection, and
  truncation maps.  Omitting any of these removes a term from the displayed
  cochain \(R_{n,M}\), so neither its first nonzero order nor its
  cohomology class is defined.
\end{proof}

\begin{rmk}[Scalar-gate boundary]\label{rmk:actual-finite-window-scalar-gate-boundary}
  In the ordinary scalar-reduced \(\lie{gl}_N\) branch, the same
  two-Hamiltonian one-cross graph gives the order-one scalar class
  \(N\,\omega(f,g)\gamma_{\mathbf 1}\), equivalently
  \(N[\bar c]\gamma_{\mathbf 1}\) in cohomology after extracting the
  \(\hbar^1\)-coefficient.  In the even block-diagonal marked habitat a
  single marker
  \(T=\lambda_0P_{\bar0}+\lambda_1P_{\bar1}\) gives
  \(N(\lambda_0-\lambda_1)\omega(f,g)\gamma_{\mathbf 1}\).  These vanish
  only under the corresponding central-character, balanced, or
  full-equivariant hypotheses; they are not non-scalar counterterms.
\end{rmk}

\begin{rmk}[Relation to regulator independence]\label{rmk:hadamard-regulator-independence}
  Theorem~\ref{thm:wt-regulator-independence-admissible} proves that,
  after the finite-window Mittag-Leffler hypotheses of
  Definition~\ref{defn:regulator-admissible-sector} are imposed, the
  Hadamard criterion is independent of the admissible exponential
  weight.  This is weight-independence among regulated coefficient
  products.  The unweighted product/direct-sum pair is separated by the
  endpoint obstruction of Theorem~\ref{thm:strict-unweighted-no-go}: the
  required coefficient Casimir has no strict product/direct-sum
  representative.
\end{rmk}
